% Options for packages loaded elsewhere
% Options for packages loaded elsewhere
\PassOptionsToPackage{unicode}{hyperref}
\PassOptionsToPackage{hyphens}{url}
\PassOptionsToPackage{dvipsnames,svgnames,x11names}{xcolor}
%
\documentclass[
  spanish,
  12pt,
]{article}
\usepackage{xcolor}
\usepackage[margin=2cm]{geometry}
\usepackage{amsmath,amssymb}
\setcounter{secnumdepth}{5}
\usepackage{iftex}
\ifPDFTeX
  \usepackage[T1]{fontenc}
  \usepackage[utf8]{inputenc}
  \usepackage{textcomp} % provide euro and other symbols
\else % if luatex or xetex
  \usepackage{unicode-math} % this also loads fontspec
  \defaultfontfeatures{Scale=MatchLowercase}
  \defaultfontfeatures[\rmfamily]{Ligatures=TeX,Scale=1}
\fi
\usepackage{lmodern}
\ifPDFTeX\else
  % xetex/luatex font selection
  \setmainfont[]{Times New Roman}
\fi
% Use upquote if available, for straight quotes in verbatim environments
\IfFileExists{upquote.sty}{\usepackage{upquote}}{}
\IfFileExists{microtype.sty}{% use microtype if available
  \usepackage[]{microtype}
  \UseMicrotypeSet[protrusion]{basicmath} % disable protrusion for tt fonts
}{}
\usepackage{setspace}
\makeatletter
\@ifundefined{KOMAClassName}{% if non-KOMA class
  \IfFileExists{parskip.sty}{%
    \usepackage{parskip}
  }{% else
    \setlength{\parindent}{0pt}
    \setlength{\parskip}{6pt plus 2pt minus 1pt}}
}{% if KOMA class
  \KOMAoptions{parskip=half}}
\makeatother
% Make \paragraph and \subparagraph free-standing
\makeatletter
\ifx\paragraph\undefined\else
  \let\oldparagraph\paragraph
  \renewcommand{\paragraph}{
    \@ifstar
      \xxxParagraphStar
      \xxxParagraphNoStar
  }
  \newcommand{\xxxParagraphStar}[1]{\oldparagraph*{#1}\mbox{}}
  \newcommand{\xxxParagraphNoStar}[1]{\oldparagraph{#1}\mbox{}}
\fi
\ifx\subparagraph\undefined\else
  \let\oldsubparagraph\subparagraph
  \renewcommand{\subparagraph}{
    \@ifstar
      \xxxSubParagraphStar
      \xxxSubParagraphNoStar
  }
  \newcommand{\xxxSubParagraphStar}[1]{\oldsubparagraph*{#1}\mbox{}}
  \newcommand{\xxxSubParagraphNoStar}[1]{\oldsubparagraph{#1}\mbox{}}
\fi
\makeatother


\usepackage{longtable,booktabs,array}
\newcounter{none} % for unnumbered tables
\usepackage{calc} % for calculating minipage widths
% Correct order of tables after \paragraph or \subparagraph
\usepackage{etoolbox}
\makeatletter
\patchcmd\longtable{\par}{\if@noskipsec\mbox{}\fi\par}{}{}
\makeatother
% Allow footnotes in longtable head/foot
\IfFileExists{footnotehyper.sty}{\usepackage{footnotehyper}}{\usepackage{footnote}}
\makesavenoteenv{longtable}
\usepackage{graphicx}
\makeatletter
\newsavebox\pandoc@box
\newcommand*\pandocbounded[1]{% scales image to fit in text height/width
  \sbox\pandoc@box{#1}%
  \Gscale@div\@tempa{\textheight}{\dimexpr\ht\pandoc@box+\dp\pandoc@box\relax}%
  \Gscale@div\@tempb{\linewidth}{\wd\pandoc@box}%
  \ifdim\@tempb\p@<\@tempa\p@\let\@tempa\@tempb\fi% select the smaller of both
  \ifdim\@tempa\p@<\p@\scalebox{\@tempa}{\usebox\pandoc@box}%
  \else\usebox{\pandoc@box}%
  \fi%
}
% Set default figure placement to htbp
\def\fps@figure{htbp}
\makeatother


% definitions for citeproc citations
\NewDocumentCommand\citeproctext{}{}
\NewDocumentCommand\citeproc{mm}{%
  \begingroup\def\citeproctext{#2}\cite{#1}\endgroup}
\makeatletter
 % allow citations to break across lines
 \let\@cite@ofmt\@firstofone
 % avoid brackets around text for \cite:
 \def\@biblabel#1{}
 \def\@cite#1#2{{#1\if@tempswa , #2\fi}}
\makeatother
\newlength{\cslhangindent}
\setlength{\cslhangindent}{1.5em}
\newlength{\csllabelwidth}
\setlength{\csllabelwidth}{3em}
\newenvironment{CSLReferences}[2] % #1 hanging-indent, #2 entry-spacing
 {\begin{list}{}{%
  \setlength{\itemindent}{0pt}
  \setlength{\leftmargin}{0pt}
  \setlength{\parsep}{0pt}
  % turn on hanging indent if param 1 is 1
  \ifodd #1
   \setlength{\leftmargin}{\cslhangindent}
   \setlength{\itemindent}{-1\cslhangindent}
  \fi
  % set entry spacing
  \setlength{\itemsep}{#2\baselineskip}}}
 {\end{list}}
\usepackage{calc}
\newcommand{\CSLBlock}[1]{\hfill\break\parbox[t]{\linewidth}{\strut\ignorespaces#1\strut}}
\newcommand{\CSLLeftMargin}[1]{\parbox[t]{\csllabelwidth}{\strut#1\strut}}
\newcommand{\CSLRightInline}[1]{\parbox[t]{\linewidth - \csllabelwidth}{\strut#1\strut}}
\newcommand{\CSLIndent}[1]{\hspace{\cslhangindent}#1}

\ifLuaTeX
\usepackage[bidi=basic,shorthands=off]{babel}
\else
\usepackage[bidi=default,shorthands=off]{babel}
\fi
\ifPDFTeX
\else
\babelfont{rm}[]{Times New Roman}
\fi
\ifLuaTeX
  \usepackage{selnolig} % disable illegal ligatures
\fi


\setlength{\emergencystretch}{3em} % prevent overfull lines

\providecommand{\tightlist}{%
  \setlength{\itemsep}{0pt}\setlength{\parskip}{0pt}}



 


\usepackage[noblocks]{authblk}
\renewcommand*{\Authsep}{, }
\renewcommand*{\Authand}{, }
\renewcommand*{\Authands}{, }
\renewcommand\Affilfont{\small}
\usepackage{etoolbox}
\AtBeginEnvironment{longtable}{\footnotesize}
\makeatletter
\@ifpackageloaded{tcolorbox}{}{\usepackage[skins,breakable]{tcolorbox}}
\@ifpackageloaded{fontawesome5}{}{\usepackage{fontawesome5}}
\definecolor{quarto-callout-color}{HTML}{909090}
\definecolor{quarto-callout-note-color}{HTML}{0758E5}
\definecolor{quarto-callout-important-color}{HTML}{CC1914}
\definecolor{quarto-callout-warning-color}{HTML}{EB9113}
\definecolor{quarto-callout-tip-color}{HTML}{00A047}
\definecolor{quarto-callout-caution-color}{HTML}{FC5300}
\definecolor{quarto-callout-color-frame}{HTML}{acacac}
\definecolor{quarto-callout-note-color-frame}{HTML}{4582ec}
\definecolor{quarto-callout-important-color-frame}{HTML}{d9534f}
\definecolor{quarto-callout-warning-color-frame}{HTML}{f0ad4e}
\definecolor{quarto-callout-tip-color-frame}{HTML}{02b875}
\definecolor{quarto-callout-caution-color-frame}{HTML}{fd7e14}
\makeatother
\makeatletter
\@ifpackageloaded{caption}{}{\usepackage{caption}}
\AtBeginDocument{%
\ifdefined\contentsname
  \renewcommand*\contentsname{Tabla de contenidos}
\else
  \newcommand\contentsname{Tabla de contenidos}
\fi
\ifdefined\listfigurename
  \renewcommand*\listfigurename{Listado de Figuras}
\else
  \newcommand\listfigurename{Listado de Figuras}
\fi
\ifdefined\listtablename
  \renewcommand*\listtablename{Listado de Tablas}
\else
  \newcommand\listtablename{Listado de Tablas}
\fi
\ifdefined\figurename
  \renewcommand*\figurename{Figura}
\else
  \newcommand\figurename{Figura}
\fi
\ifdefined\tablename
  \renewcommand*\tablename{Tabla}
\else
  \newcommand\tablename{Tabla}
\fi
}
\@ifpackageloaded{float}{}{\usepackage{float}}
\floatstyle{ruled}
\@ifundefined{c@chapter}{\newfloat{codelisting}{h}{lop}}{\newfloat{codelisting}{h}{lop}[chapter]}
\floatname{codelisting}{Listado}
\newcommand*\listoflistings{\listof{codelisting}{Listado de Listados}}
\makeatother
\makeatletter
\makeatother
\makeatletter
\@ifpackageloaded{caption}{}{\usepackage{caption}}
\@ifpackageloaded{subcaption}{}{\usepackage{subcaption}}
\makeatother
\usepackage{bookmark}
\IfFileExists{xurl.sty}{\usepackage{xurl}}{} % add URL line breaks if available
\urlstyle{same}
\hypersetup{
  pdftitle={{[}Working title{]} Percepciones del cambio tecnológico y preferencias de política social de los trabajadores en 27 países de la OCDE},
  pdfauthor={Kevin Carrasco},
  pdflang={es},
  colorlinks=true,
  linkcolor={blue},
  filecolor={Maroon},
  citecolor={Blue},
  urlcolor={Blue},
  pdfcreator={LaTeX via pandoc}}


\title{{[}Working title{]} Percepciones del cambio tecnológico y
preferencias de política social de los trabajadores en 27 países de la
OCDE}


  \author{Kevin Carrasco}
            \affil{%
                  Núcleo Milenio sobre Evolución del Trabajo (M-NEW)
              }
      
\date{}
\begin{document}
\maketitle


\setstretch{1.15}
\section{Introducción}\label{introducciuxf3n}

La difusión de la inteligencia artificial (IA), y en particular de la IA
generativa, ha reabierto la pregunta sobre cómo deben responder los
Estados a las transformaciones tecnológicas del trabajo. Las olas
previas de automatización afectaron sobre todo a tareas rutinarias y a
ocupaciones de calificación baja y media
(\citeproc{ref-acemoglu_tasks_2022}{Acemoglu \& Restrepo, 2022};
\citeproc{ref-autor_why_2015}{Autor, 2015};
\citeproc{ref-frey_future_2017}{Frey \& Osborne, 2017}). Las
estimaciones recientes indican que la exposición a la IA generativa
alcanza también a ocupaciones cognitivas de mayor calificación y
salario, como las de las finanzas, el derecho o la medicina, y que en
muchas de ellas es más probable que transforme las tareas a que elimine
los puestos de trabajo (\citeproc{ref-eloundou_gpts_2024}{Eloundou
et~al., 2024}; \citeproc{ref-gmyrek_generative_2023}{Gmyrek et~al.,
2023}; \citeproc{ref-oecd_employment_2023}{OECD, 2023}). Esta ampliación
de la exposición vuelve más incierto quiénes ganarán y quiénes perderán
con el cambio tecnológico
(\citeproc{ref-haslberger_rage_2025}{Haslberger et~al., 2025}) y abre
una disputa sobre qué debe hacer el Estado frente a él. Algunas
propuestas apuntan a facilitar la adaptación mediante la inversión en
educación, reentrenamiento e infraestructura digital, otras a compensar
las pérdidas mediante prestaciones más generosas o una renta básica, y
otras a proteger los empleos y las industrias existentes mediante
impuestos a la automatización, subsidios a las industrias afectadas, la
preferencia por los trabajadores del país o los aranceles
(\citeproc{ref-busemeyer_preferred_2023}{Busemeyer et~al., 2023};
\citeproc{ref-colantone_trade_2018}{Colantone \& Stanig, 2018};
\citeproc{ref-gallego_technological_2022}{Gallego et~al., 2022};
\citeproc{ref-lee_luddite_2024}{Lee, 2024}).

Este artículo estudia las preferencias de la ciudadanía sobre estas
respuestas desde la sociología del Estado de bienestar y la economía
política comparada. Su objeto son las preferencias de política social
frente al cambio tecnológico, que comprenden tanto la cantidad de
intervención estatal que las personas demandan como el tipo de respuesta
que priorizan. Los factores asociados que examina son las percepciones
sobre el cambio tecnológico, referidas al propio empleo y a los efectos
de la IA sobre la sociedad, y el contexto institucional en que esas
percepciones se traducen en preferencias.

La investigación sobre las consecuencias políticas del cambio
tecnológico ha mostrado que el riesgo de automatización se asocia con
las preferencias de política social, aunque los mecanismos que conectan
ambos fenómenos siguen siendo poco comprendidos
(\citeproc{ref-gallego_automation_2022}{Gallego \& Kurer, 2022}).
Quienes trabajan en ocupaciones más expuestas o perciben un mayor riesgo
de ser reemplazados tienden a apoyar más la redistribución y la
compensación, sin un aumento equivalente del apoyo a la inversión social
(\citeproc{ref-busemeyer_social_2022}{Busemeyer \& Sahm, 2022};
\citeproc{ref-busemeyer_dealing_2023}{Busemeyer \& Tober, 2023};
\citeproc{ref-kurer_automation_2022}{Kurer \& Häusermann, 2022};
\citeproc{ref-thewissen_automation_2019}{Thewissen \& Rueda, 2019}).
Otros estudios encuentran que la preocupación por el desplazamiento se
orienta hacia medidas que frenan el cambio tecnológico y apenas modifica
la demanda de compensación (\citeproc{ref-bicchi_unpacking_2025}{Bicchi
et~al., 2025}; \citeproc{ref-gallego_technological_2022}{Gallego et~al.,
2022}). Una tercera línea muestra que los trabajadores amenazados por la
automatización expresan mayor hostilidad hacia el comercio y la
inmigración y mayor apoyo a la derecha radical
(\citeproc{ref-anelli_individual_2021}{Anelli et~al., 2021};
\citeproc{ref-im_losers_2019}{Im et~al., 2019};
\citeproc{ref-wu_misattributed_2022}{Wu, 2022}). Estos resultados
ofrecen respuestas distintas a la pregunta de hacia dónde se dirige la
demanda de quienes perciben el cambio tecnológico como una amenaza.

Hasta ahora, esta literatura comparte tres rasgos que dificultan
responder esa pregunta. El primero es el supuesto de que el vínculo
entre tecnología y preferencias pasa por el interés propio, de modo que
la percepción relevante es la que cada persona tiene sobre su propio
empleo. La investigación sobre actitudes redistributivas muestra, sin
embargo, que el apoyo a las políticas sociales responde también a normas
compartidas de justicia (\citeproc{ref-cavaille_fair_2025}{Cavaillé,
2025}), y la tecnología genera expectativas sobre lo que ocurrirá en la
sociedad que pueden activar esas normas. El segundo rasgo es que las
respuestas se han estudiado una por una o mediante asignaciones
presupuestarias que fijan de antemano el monto total
(\citeproc{ref-busemeyer_dealing_2023}{Busemeyer \& Tober, 2023}), lo
que impide saber si el cambio tecnológico aumenta la demanda de
intervención estatal o la reorienta hacia respuestas específicas. El
tercero es que las medidas que resguardan el mercado laboral nacional de
la competencia externa rara vez se han analizado como respuesta al
cambio tecnológico, junto a la inversión, la compensación y las medidas
que frenan la tecnología.

Este artículo sostiene que las preferencias frente al cambio tecnológico
tienen dos componentes que responden a lógicas distintas. El volumen
corresponde a la cantidad de intervención estatal que se demanda,
mientras que la dirección refiere a la respuesta que se prioriza entre
la inversión, la compensación y la protección. Las percepciones sobre el
cambio tecnológico pueden reorientar la demanda sin aumentarla, y el
sentido de esa reorientación depende del motivo que activan, sea el
interés propio o la adhesión a normas de justicia. A partir de la
encuesta Risks That Matter de la OCDE de 2024, que cubre a más de 27.000
personas de entre 18 y 64 años en 27 países
(\citeproc{ref-oecd_risks_2024}{OECD, 2024};
\citeproc{ref-thomas_survey_2025}{Thomas et~al., 2025}), el artículo
contribuye a la investigación sobre cambio tecnológico y Estado de
bienestar al separar el volumen de la dirección de la demanda, al
incorporar a la protección las medidas que resguardan el mercado
nacional de la competencia externa y al evaluar el papel de
instituciones que corresponden a cada repertorio. La pregunta que
orienta el estudio es cómo se relacionan las percepciones sobre el
cambio tecnológico con la cantidad de intervención estatal que demandan
las personas y con el tipo de respuesta que priorizan, y cómo condiciona
el contexto institucional esa relación.

Las percepciones sobre el propio empleo activan el interés propio y
orientan la demanda hacia las respuestas que protegen o amplían la
posición de cada persona, en una dirección que depende de si la amenaza
se atribuye a riesgos que el Estado debe asegurar, a la tecnología y las
empresas que la adoptan o a la competencia extranjera. Las expectativas
sobre los efectos sociales de la IA activan, por su parte, normas de
justicia que definen si el cambio tecnológico se percibe como un
problema de adaptación o como un problema distributivo, y ambas
percepciones se traducen en preferencias dentro de instituciones que
pueden volver más disponible el repertorio correspondiente o reducir la
demanda de expandirlo.

De este modo, el estudio relaciona tres conjuntos de conceptos. Las
preferencias frente al cambio tecnológico, en su volumen y su dirección,
constituyen el objeto de estudio. Las percepciones sobre el propio
empleo y las expectativas sobre los efectos sociales de la IA
constituyen sus factores individuales, y las prestaciones por desempleo
y las políticas de mercado laboral constituyen su contexto. Las
secciones siguientes precisan cada concepto y presentan las hipótesis
que se derivan de ellos.

\subsection{Preferencias frente al cambio
tecnológico}\label{preferencias-frente-al-cambio-tecnoluxf3gico}

Las preferencias frente al cambio tecnológico se refieren al apoyo de
las personas a las respuestas que el Estado puede adoptar ante la
transformación tecnológica del trabajo. Estas respuestas pueden
distinguirse según el momento en que intervienen y según quién asume los
costos del ajuste. La inversión interviene antes de que los riesgos se
materialicen y busca que las personas y la economía puedan aprovechar
las nuevas tecnologías mediante la educación, el reentrenamiento y la
infraestructura digital. Se vincula con la perspectiva de la inversión
social, que orienta la política social hacia la formación de capacidades
a lo largo de la vida
(\citeproc{ref-garritzmann_public_2018}{Garritzmann et~al., 2018};
\citeproc{ref-morel_towards_2012}{Morel et~al., 2012}). La compensación
interviene una vez que los riesgos se materializan, mediante
transferencias que aseguran el ingreso de quienes pierden su empleo o
ven reducidas sus oportunidades, y responde a la lógica del
aseguramiento social (\citeproc{ref-iversen_asset_2001}{Iversen \&
Soskice, 2001}; \citeproc{ref-rehm_risks_2009}{Rehm, 2009}). Su
discusión se ha renovado con la expansión del trabajo mediado por
plataformas, que externaliza responsabilidades propias de la relación
laboral (\citeproc{ref-vallas_what_2020}{Vallas \& Schor, 2020}) y
presenta brechas de cobertura que obligan a adaptar los sistemas de
protección social (\citeproc{ref-behrendt_social_2019}{Behrendt et~al.,
2019}).

La protección busca evitar que los costos del cambio recaigan sobre los
trabajadores de los sectores afectados, preservando los empleos y las
industrias existentes frente a la sustitución por máquinas o por
competidores extranjeros. A diferencia de la protección social, que
asegura el ingreso y en este estudio corresponde a la compensación, su
objeto son los puestos de trabajo y las industrias. Comprende medidas
que frenan o encarecen la sustitución tecnológica, como los impuestos a
la automatización y los subsidios a los sectores más afectados, y
medidas que resguardan el mercado laboral nacional de la competencia
externa, como la preferencia por los trabajadores del país en la
contratación y los aranceles. Las primeras corresponden a lo que la
literatura ha denominado proteccionismo tecnológico
(\citeproc{ref-bicchi_unpacking_2025}{Bicchi et~al., 2025};
\citeproc{ref-gallego_technological_2022}{Gallego et~al., 2022}) y
trasladan parte de los costos del ajuste hacia las empresas que
automatizan o hacia el presupuesto público. Las segundas se vinculan con
el nacionalismo económico, que ha ganado apoyo en las regiones más
expuestas a la competencia de las importaciones
(\citeproc{ref-colantone_trade_2018}{Colantone \& Stanig, 2018}), y con
criterios de merecimiento que, en los países europeos, sitúan a los
inmigrantes en el último lugar entre los grupos que merecen apoyo
(\citeproc{ref-vanoorschot_making_2006}{Van Oorschot, 2006}). Ambos
tipos de medidas difieren en el agente al que apuntan y comparten el
objetivo de preservar las posiciones existentes frente al ajuste. Ese
objetivo distingue a la protección de la inversión, que prepara a las
personas para el cambio, y de la compensación, que atenúa las pérdidas
que este produce.

Los tres repertorios expresan intereses distintos en el mercado laboral.
En mercados laborales dualizados, quienes cuentan con empleos estables
priorizan la protección del empleo, mientras que quienes tienen
inserciones más precarias priorizan las políticas activas y pasivas
dirigidas a los desempleados
(\citeproc{ref-emmenegger_age_2012}{Emmenegger et~al., 2012};
\citeproc{ref-rueda_social_2007}{Rueda, 2007}). La protección del
empleo, la protección frente al desempleo y la formación de
calificaciones funcionan, además, como complementos institucionales
(\citeproc{ref-estevez-abe_social_2001}{Estevez-Abe et~al., 2001}). Las
preferencias entre repertorios reflejan, por lo tanto, un conflicto
distributivo sobre quién asume los costos del cambio tecnológico.

Estas preferencias tienen dos componentes. El volumen corresponde al
nivel general de apoyo a las respuestas frente al cambio tecnológico y
expresa cuánta intervención estatal se demanda. La dirección corresponde
al apoyo a cada repertorio en relación con ese nivel general y expresa
qué respuesta se prioriza. Dos personas pueden diferir en el volumen y
coincidir en la dirección, como ocurre cuando una apoya con intensidad
todas las respuestas y otra las apoya moderadamente, pero ambas
privilegian la compensación por sobre la inversión. La distinción retoma
la diferencia entre el tamaño y el contenido de la intervención estatal,
ya que los conflictos distributivos sobre qué políticas expandir siguen
dinámicas distintas de los debates sobre el tamaño general del Estado de
bienestar (\citeproc{ref-busemeyer_public_2017}{Busemeyer \&
Garritzmann, 2017}; \citeproc{ref-hausermann_politics_2019}{Häusermann
et~al., 2019}).

La evidencia empírica muestra que esta distinción es necesaria. Cuando
las medidas se evalúan una por una, sin restricciones presupuestarias,
el apoyo tiende a ser alto para casi todas y disminuye considerablemente
al introducir costos o disyuntivas explícitas
(\citeproc{ref-busemeyer_public_2017}{Busemeyer \& Garritzmann, 2017};
\citeproc{ref-hausermann_politics_2019}{Häusermann et~al., 2019}). La
inversión social, en particular la educación, concentra un apoyo amplio
(\citeproc{ref-busemeyer_public_2017}{Busemeyer \& Garritzmann, 2017}),
y entre las respuestas al cambio tecnológico la inversión es la que
recibe mayor apoyo (\citeproc{ref-churchill_policy_2025}{Churchill
et~al., 2025}). Busemeyer \& Tober
(\citeproc{ref-busemeyer_dealing_2023}{2023}) abordan la disyuntiva
entre respuestas mediante una asignación presupuestaria que obliga a
distribuir recursos entre opciones. Ese diseño identifica prioridades,
pero fija de antemano el monto total, por lo que no permite saber si el
cambio tecnológico aumenta la demanda de intervención estatal o solo la
reorienta. Separar el volumen de la dirección permite responder ambas
preguntas a partir de las mismas preferencias.

\subsection{Amenaza y beneficio tecnológico en el propio
empleo}\label{amenaza-y-beneficio-tecnoluxf3gico-en-el-propio-empleo}

La percepción de amenaza tecnológica se refiere a la probabilidad que
cada persona asigna a perder su empleo por la automatización, por la
competencia mediada por plataformas o por la obsolescencia de sus
habilidades. El beneficio esperado se refiere a la probabilidad de que
la tecnología mejore las condiciones del propio trabajo. Ambas
percepciones constituyen dimensiones distintas, porque la tecnología
sustituye trabajo humano en algunas tareas y lo complementa en otras
(\citeproc{ref-autor_why_2015}{Autor, 2015}). En España, por ejemplo, la
mayoría de los trabajadores evalúa positivamente el efecto de las nuevas
tecnologías en su lugar de trabajo, aunque existe una minoría preocupada
por el desplazamiento (\citeproc{ref-gallego_technological_2022}{Gallego
et~al., 2022}). La amenaza percibida se distribuye de manera desigual
entre ocupaciones y grupos sociales, y sigue en buena medida la
exposición objetiva de las tareas
(\citeproc{ref-dekker_fear_2017}{Dekker et~al., 2017};
\citeproc{ref-gallego_technological_2022}{Gallego et~al., 2022}). La
percepción de amenaza ubica así a cada persona en una posición de riesgo
frente al cambio tecnológico, del mismo modo en que el riesgo de empleo
atípico o de desempleo permite distinguir a quienes están dentro y fuera
de la protección del mercado laboral
(\citeproc{ref-schwander_who_2013}{Schwander \& Häusermann, 2013}).

Respecto del volumen de la demanda, la teoría más establecida sostiene
que el Estado de bienestar funciona como un seguro, de modo que quienes
enfrentan mayores riesgos en el mercado laboral demandan más protección
social (\citeproc{ref-cusack_risks_2006}{Cusack et~al., 2006};
\citeproc{ref-iversen_asset_2001}{Iversen \& Soskice, 2001};
\citeproc{ref-rehm_risks_2009}{Rehm, 2009}). La inseguridad laboral
subjetiva aumenta la demanda de redistribución, en especial entre
quienes temen un desempleo prolongado
(\citeproc{ref-marx_effect_2014}{Marx, 2014}), y el empleo en
ocupaciones intensivas en tareas rutinarias se asocia con un mayor apoyo
al seguro público frente a la pérdida de ingresos
(\citeproc{ref-thewissen_automation_2019}{Thewissen \& Rueda, 2019}).
Desde esta perspectiva, la amenaza tecnológica debería aumentar la
cantidad de intervención estatal que demandan quienes se sienten
amenazados.

Hay razones, sin embargo, para dudar de que esa traducción sea directa.
El riesgo de automatización se asocia con más apoyo a la redistribución
y con menos apoyo a la inversión social, sin relación con el apoyo a una
renta básica (\citeproc{ref-busemeyer_social_2022}{Busemeyer \& Sahm,
2022}), y los trabajadores que se sienten amenazados apoyan la
compensación pasiva sin aumentar su apoyo a las políticas activas
(\citeproc{ref-kurer_automation_2022}{Kurer \& Häusermann, 2022}). La
disposición a ampliar la intervención estatal depende, además, de la
confianza en que el Estado actuará con eficiencia e imparcialidad
(\citeproc{ref-rothstein_just_1998}{Rothstein, 1998};
\citeproc{ref-svallfors_government_2013}{Stefan Svallfors, 2013}). Si
quienes se sienten amenazados dudan de la capacidad estatal para
gestionar el cambio tecnológico, su demanda puede concentrarse en
respuestas específicas sin aumentar en su conjunto. La lógica del
aseguramiento permite, con todo, formular una primera hipótesis:

\begin{tcolorbox}[enhanced jigsaw, arc=.35mm, bottomrule=.15mm, breakable, colback=white, colframe=quarto-callout-note-color-frame, left=2mm, leftrule=.75mm, opacityback=0, rightrule=.15mm, toprule=.15mm]

H1: Quienes perciben una mayor amenaza tecnológica sobre su empleo
demandarán más intervención estatal frente al cambio tecnológico.

\end{tcolorbox}

Respecto de la dirección, la respuesta preferida depende de a qué se
atribuye la amenaza. La demanda de protección comercial tras la pérdida
de un empleo varía según la causa que se le atribuye y es mayor cuando
esta se asocia al comercio internacional
(\citeproc{ref-ditella_labour_2020}{\textbf{ditella\_labour\_2020?}}).
En el ámbito tecnológico, la preocupación por la descalificación se
asocia con demandas de gravar a las empresas que se consideran
responsables de los cambios
(\citeproc{ref-churchill_policy_2025}{Churchill et~al., 2025}). A partir
de la atribución pueden formularse dos tesis sobre hacia dónde se dirige
la amenaza tecnológica.

La tesis de la compensación sostiene que la amenaza tecnológica se
canaliza hacia el Estado como asegurador. Las prestaciones aseguran el
ingreso de manera inmediata si el riesgo se materializa, mientras que
los beneficios de la formación son inciertos, llegan a mediano plazo y
dependen de la capacidad de cada persona para reconvertirse
(\citeproc{ref-busemeyer_dealing_2023}{Busemeyer \& Tober, 2023}). Los
estudios que usan medidas subjetivas u ocupacionales del riesgo de
automatización encuentran, entre quienes enfrentan más riesgo, un mayor
apoyo a la compensación y un apoyo igual o menor a la inversión social
(\citeproc{ref-busemeyer_social_2022}{Busemeyer \& Sahm, 2022};
\citeproc{ref-busemeyer_dealing_2023}{Busemeyer \& Tober, 2023};
\citeproc{ref-kurer_automation_2022}{Kurer \& Häusermann, 2022}).

\begin{tcolorbox}[enhanced jigsaw, arc=.35mm, bottomrule=.15mm, breakable, colback=white, colframe=quarto-callout-note-color-frame, left=2mm, leftrule=.75mm, opacityback=0, rightrule=.15mm, toprule=.15mm]

H2a: Quienes perciben una mayor amenaza tecnológica sobre su empleo
darán mayor prioridad a la compensación y menor prioridad a la
inversión.

\end{tcolorbox}

La tesis de la protección sostiene que quienes se sienten amenazados
buscan preservar su posición frente al ajuste, y que esa demanda puede
seguir dos vías. La primera dirige la amenaza contra su fuente. Para
quien teme perder su empleo, las medidas que frenan la automatización
actúan sobre el origen del riesgo, mientras que la compensación solo
atenúa sus consecuencias. Por un lado, en España, los trabajadores
preocupados por el desplazamiento prefieren políticas que frenen el
cambio tecnológico y, en menor medida, demandan más compensación
(\citeproc{ref-bicchi_unpacking_2025}{Bicchi et~al., 2025};
\citeproc{ref-gallego_technological_2022}{Gallego et~al., 2022}). Por
otro lado, en el Reino Unido y Suecia la preocupación por la
automatización se asocia con el apoyo a restringirla, sobre todo cuando
las salvaguardas laborales se perciben débiles
(\citeproc{ref-lee_luddite_2024}{Lee, 2024}).

La segunda vía desplaza la amenaza tecnológica hacia la competencia
extranjera. Los trabajadores más expuestos a la automatización muestran
mayor oposición al comercio y a la inmigración, lo que sugiere que
atribuyen a la apertura pérdidas que produce la tecnología
(\citeproc{ref-wu_misattributed_2022}{Wu, 2022}), y apoyan en mayor
medida a partidos de derecha radical que combinan la preferencia
nacional con el proteccionismo económico
(\citeproc{ref-anelli_individual_2021}{Anelli et~al., 2021};
\citeproc{ref-im_losers_2019}{Im et~al., 2019};
\citeproc{ref-kurer_declining_2020}{Kurer, 2020}). La evidencia
experimental indica, con todo, que la sustitución por robots genera
menos demanda de medidas preventivas que la deslocalización o la
contratación de trabajadores extranjeros
(\citeproc{ref-kuo_automation_2024}{Kuo et~al., 2024}). Por cualquiera
de las dos vías, la amenaza tecnológica debería aumentar la prioridad de
la protección.

\begin{tcolorbox}[enhanced jigsaw, arc=.35mm, bottomrule=.15mm, breakable, colback=white, colframe=quarto-callout-note-color-frame, left=2mm, leftrule=.75mm, opacityback=0, rightrule=.15mm, toprule=.15mm]

H2b: Quienes perciben una mayor amenaza tecnológica sobre su empleo
darán mayor prioridad a la protección.

\end{tcolorbox}

El beneficio esperado de la tecnología sitúa a la persona en una
posición distinta. Quien anticipa que la tecnología mejorará sus
condiciones de trabajo ocupa una posición complementaria con el cambio
tecnológico, más cercana a la ampliación de sus tareas que a su
sustitución. Según Gmyrek et~al.
(\citeproc{ref-gmyrek_generative_2023}{2023}), esa ampliación es el
efecto más probable de la IA generativa en muchas ocupaciones. La
inversión en educación, reentrenamiento e infraestructura amplía las
oportunidades de aprovechar la tecnología, mientras que las medidas de
protección, al frenar el ajuste, ponen en riesgo esos beneficios
(\citeproc{ref-gmyrek_generative_2023}{Gmyrek et~al., 2023}). Busemeyer
et~al. (\citeproc{ref-busemeyer_preferred_2023}{2023}) esperan que
quienes usan tecnología de forma constante en el trabajo, tienen
educación universitaria y mayores ingresos prefieran la inversión, y
encuentran que incluso en esos grupos el aumento del riesgo percibido
desplaza el apoyo hacia las políticas compensatorias. A su vez, la
información de que las ganancias de la automatización se distribuyen
ampliamente aumenta el apoyo a acelerarla, en especial entre los
trabajadores menos expuestos (\citeproc{ref-lee_luddite_2024}{Lee,
2024}).

\begin{tcolorbox}[enhanced jigsaw, arc=.35mm, bottomrule=.15mm, breakable, colback=white, colframe=quarto-callout-note-color-frame, left=2mm, leftrule=.75mm, opacityback=0, rightrule=.15mm, toprule=.15mm]

H3: Quienes esperan mayores beneficios de la tecnología en su empleo
darán mayor prioridad a la inversión y menor prioridad a la protección.

\end{tcolorbox}

\subsection{Expectativas sobre los efectos sociales de la inteligencia
artificial}\label{expectativas-sobre-los-efectos-sociales-de-la-inteligencia-artificial}

Las expectativas sobre los efectos sociales de la IA se refieren a lo
que las personas anticipan que ocurrirá en la sociedad, como un aumento
de la desigualdad y del desempleo o la creación de empleos y la mejora
de las condiciones de trabajo. Estas expectativas funcionan como
diagnósticos sobre el tipo de transformación que producirá la
tecnología. Si la IA se percibe como fuente de desigualdad y desempleo,
el cambio tecnológico aparece como un problema distributivo en el que
unos concentran las ganancias y otros asumen los costos. Si se percibe
como fuente de empleo y de mejores condiciones de trabajo, aparece como
un problema de adaptación.

Estos diagnósticos se vinculan con las preferencias a través de normas
de justicia. La investigación sobre actitudes redistributivas identifica
dos motivos en el apoyo a las políticas sociales, que son el interés en
el propio ingreso esperado y la adhesión a normas compartidas de
justicia, y muestra que el segundo tiende a predominar salvo cuando las
consecuencias de una política para el propio ingreso son grandes y
ciertas (\citeproc{ref-cavaille_fair_2025}{Cavaillé, 2025}). La economía
moral del Estado de bienestar descansa en normas de reciprocidad y en
juicios sobre una distribución justa de cargas y beneficios
(\citeproc{ref-mau_moral_2014}{Mau, 2014}), y un cambio que concentra
las ganancias en quienes adoptan la tecnología puede percibirse como una
vulneración de esas normas. El cambio tecnológico ofrece condiciones
para que este motivo tenga peso, porque solo una fracción de los
trabajadores percibe un riesgo alto para su propio empleo
(\citeproc{ref-kurer_automation_2022}{Kurer \& Häusermann, 2022}).

La evidencia reciente es consistente con este razonamiento. La
exposición directa a la IA generativa aumenta el apoyo a políticas
sociales activadoras sin elevar el riesgo percibido, un resultado que
sus autores atribuyen a la preocupación por quienes podrían verse
perjudicados (\citeproc{ref-haslberger_rage_2025}{Haslberger et~al.,
2025}). A su vez, la preocupación general por los efectos de la
tecnología se asocia con el apoyo al proteccionismo tecnológico,
mientras que el riesgo de sustitución se relaciona solo modestamente con
algunas políticas redistributivas
(\citeproc{ref-bicchi_unpacking_2025}{Bicchi et~al., 2025}).

Respecto del volumen, un diagnóstico distributivo convierte el cambio
tecnológico en un problema colectivo que reclama la acción del Estado. A
nivel agregado, la desigualdad se asocia con una mayor demanda de
redistribución (\citeproc{ref-finseraas_income_2009}{Finseraas, 2009}),
y la exposición de las economías a riesgos externos ha ido acompañada de
una expansión del gasto social, que la hipótesis de compensación
interpreta como una demanda de aseguramiento frente a esos riesgos
(\citeproc{ref-rodrik_why_1998}{Rodrik, 1998};
\citeproc{ref-walter_globalization_2010}{Walter, 2010}).

\begin{tcolorbox}[enhanced jigsaw, arc=.35mm, bottomrule=.15mm, breakable, colback=white, colframe=quarto-callout-note-color-frame, left=2mm, leftrule=.75mm, opacityback=0, rightrule=.15mm, toprule=.15mm]

H4: Las expectativas negativas sobre los efectos sociales de la IA
aumentarán la demanda de intervención estatal frente al cambio
tecnológico.

\end{tcolorbox}

Respecto de la dirección, un diagnóstico distributivo favorece las
respuestas que corrigen la distribución de los costos. La compensación
redistribuye hacia quienes pierden, y la protección traslada los costos
del ajuste desde los trabajadores de los sectores afectados hacia las
empresas que automatizan, el presupuesto público o la competencia
extranjera. Un diagnóstico de adaptación favorece, en cambio, la
inversión, porque sitúa el problema en la preparación de las personas y
del país para aprovechar las oportunidades esperadas y reduce la
urgencia de compensar pérdidas o de frenar el ajuste.

\begin{tcolorbox}[enhanced jigsaw, arc=.35mm, bottomrule=.15mm, breakable, colback=white, colframe=quarto-callout-note-color-frame, left=2mm, leftrule=.75mm, opacityback=0, rightrule=.15mm, toprule=.15mm]

H5a: Las expectativas negativas sobre los efectos sociales de la IA se
asociarán con una mayor prioridad de la compensación y de la protección.

H5b: Las expectativas positivas sobre los efectos sociales de la IA se
asociarán con una mayor prioridad de la inversión y una menor prioridad
de la compensación y de la protección.

\end{tcolorbox}

\subsection{Contexto institucional}\label{contexto-institucional}

Las preferencias de política social se forman en contextos
institucionales que contribuyen a moldearlas. La investigación comparada
ha mostrado que el apoyo a la redistribución varía entre regímenes de
bienestar (\citeproc{ref-esping-andersen_three_1993}{Esping-Andersen,
1990}; \citeproc{ref-jaeger_welfare_2006}{Jæger, 2006};
\citeproc{ref-svallfors_worlds_1997}{S. Svallfors, 1997}), que las
instituciones de cobertura amplia generan coaliciones de apoyo más
extensas que las focalizadas (\citeproc{ref-korpi_paradox_1998}{Korpi \&
Palme, 1998}) y que el diseño institucional influye en el apoyo a las
políticas a través de los juicios sobre su justicia y sobre el
merecimiento de sus beneficiarios
(\citeproc{ref-larsen_institutional_2008}{Larsen, 2008};
\citeproc{ref-rothstein_just_1998}{Rothstein, 1998}).

Las tipologías de regímenes describen con dificultad a los Estados de
bienestar reales, que rara vez se ajustan a un tipo puro
(\citeproc{ref-arts_three_2002}{Arts \& Gelissen, 2002}), y los efectos
de la protección social sobre la experiencia del riesgo dependen de
componentes específicos. Las distintas formas de inseguridad laboral
responden a instrumentos de protección distintos, mientras que la
generosidad general del Estado de bienestar no tiene un efecto
sistemático sobre ellas (\citeproc{ref-anderson_workers_2007}{Anderson
\& Pontusson, 2007}). Las políticas que aseguran el ingreso y la
empleabilidad reducen la inseguridad laboral percibida
(\citeproc{ref-chung_institutions_2011}{Chung \& Van Oorschot, 2011}), y
el temor a los robots es, en alguna medida, mayor donde los trabajadores
están menos protegidos frente al mercado
(\citeproc{ref-dekker_fear_2017}{Dekker et~al., 2017}). En el caso de la
automatización, el tamaño del Estado de bienestar modera la relación
entre el riesgo y las preferencias de política social
(\citeproc{ref-busemeyer_social_2022}{Busemeyer \& Sahm, 2022}).

Estos resultados justifican desagregar el contexto institucional según
los repertorios. Las prestaciones por desempleo corresponden a la
compensación y las políticas de mercado laboral, a la inversión. En este
estudio, las políticas de mercado laboral se refieren a las políticas
activas, que buscan que las personas consigan o conserven un empleo
mediante la capacitación, el reentrenamiento y otros programas. Esta
correspondencia sigue la lógica de complementariedad institucional entre
la protección frente al desempleo y la formación de calificaciones
(\citeproc{ref-estevez-abe_social_2001}{Estevez-Abe et~al., 2001}).

La literatura sobre retroalimentación de las políticas públicas sostiene
que las instituciones existentes moldean las preferencias porque
distribuyen recursos y ofrecen marcos para interpretar los problemas
(\citeproc{ref-pierson_when_1993}{Pierson, 1993}). Una implicancia de
este argumento es que las instituciones vuelven disponibles ciertas
respuestas. Donde un repertorio está institucionalizado, las personas lo
conocen, lo consideran viable y lo asocian con la manera habitual de
enfrentar los riesgos, de modo que quienes perciben amenazas tenderían a
canalizar su demanda hacia él. A este mecanismo puede sumarse otro de
carácter relacional. Las instituciones restringen las prácticas de las
personas, esas prácticas se interpretan como expresión de preferencias
compartidas y las respuestas institucionalizadas adquieren así una
validez social aparente (\citeproc{ref-puga_ideology_2023}{Puga \& Moya,
2023}). En China, el riesgo de automatización y el apoyo a la
responsabilidad estatal se refuerzan mutuamente donde los sistemas
locales de bienestar son más generosos (\citeproc{ref-fan_how_2024}{Fan
et~al., 2024}), un patrón consistente con este argumento. La expectativa
de disponibilidad compite con el argumento de que las instituciones que
uniforman los riesgos, como una protección del empleo fuerte o
prestaciones desvinculadas de la participación laboral, reducen el peso
del riesgo individual en las preferencias
(\citeproc{ref-gingrich_preferences_2012}{Gingrich \& Ansell, 2012}).

\begin{tcolorbox}[enhanced jigsaw, arc=.35mm, bottomrule=.15mm, breakable, colback=white, colframe=quarto-callout-note-color-frame, left=2mm, leftrule=.75mm, opacityback=0, rightrule=.15mm, toprule=.15mm]

H6a: En los países con mayor gasto en políticas de mercado laboral, la
amenaza tecnológica y las expectativas negativas sobre la IA se
traducirán en mayor medida en prioridad para la inversión, y en los
países con mayor cobertura de las prestaciones por desempleo, en
prioridad para la compensación.

\end{tcolorbox}

El modelo termostático plantea una expectativa diferente, referida al
nivel del apoyo. Según este modelo, la opinión pública ajusta su demanda
de más intervención en un ámbito en función de lo que el Estado ya hace
en él, de modo que la demanda relativa disminuye donde la política es
más generosa (\citeproc{ref-soroka_degrees_2010}{Soroka \& Wlezien,
2010}; \citeproc{ref-wlezien_public_1995a}{Wlezien, 1995}). En el ámbito
tecnológico, Busemeyer \& Tober
(\citeproc{ref-busemeyer_dealing_2023}{2023}) encuentran que el apoyo a
las prestaciones es algo menor en los países con tasas de reemplazo más
altas, lo que interpretan como un efecto de retroalimentación que limita
su propia expansión. Ambas expectativas pueden operar a la vez, ya que
en esos mismos datos las diferencias entre contextos se concentran entre
las personas de bajo riesgo y se reducen a medida que el riesgo
percibido aumenta.

\begin{tcolorbox}[enhanced jigsaw, arc=.35mm, bottomrule=.15mm, breakable, colback=white, colframe=quarto-callout-note-color-frame, left=2mm, leftrule=.75mm, opacityback=0, rightrule=.15mm, toprule=.15mm]

H6b: En los países con mayor gasto en políticas de mercado laboral, la
inversión tendrá menor prioridad, y en los países con mayor cobertura de
las prestaciones por desempleo, la compensación tendrá menor prioridad.

\end{tcolorbox}

\section{Datos, variables y métodos}\label{datos-variables-y-muxe9todos}

\subsection{Datos}\label{datos}

Este estudio utiliza datos de la ola más reciente (2024) de la encuesta
Risks That Matter (RTM) de la OCDE (\citeproc{ref-oecd_risks_2024}{OECD,
2024}), una encuesta comparada de opinión pública que cubre a 27.229
personas de entre 18 y 64 años en 27 países miembros (Alemania, Austria,
Bélgica, Canadá, Chile, Corea, Dinamarca, Eslovenia, España, Estados
Unidos, Estonia, Finlandia, Francia, Grecia, Irlanda, Israel, Italia,
Letonia, Lituania, México, Noruega, Países Bajos, Polonia, Portugal,
Reino Unido, Suiza y Türkiye). La encuesta se aplica en línea y cuenta
con ponderadores de postestratificación calibrados por edad, sexo,
región y nivel educacional (\citeproc{ref-thomas_survey_2025}{Thomas
et~al., 2025}). La ola de 2024 incorpora preguntas sobre los efectos
esperados de la inteligencia artificial en el mercado laboral y sobre el
apoyo a distintas respuestas públicas frente a la digitalización y el
cambio tecnológico.

El análisis se restringe a las personas que en la semana previa a la
entrevista tenían un trabajo remunerado o estaban temporalmente ausentes
de él por razones de salud u otras (n = 22.063). Esta restricción sigue
la estructura del cuestionario, ya que las preguntas sobre la amenaza y
el beneficio de la tecnología para el propio empleo solo se aplican a
quienes trabajan, de modo que las 5.166 personas no ocupadas, busquen o
no empleo, quedan fuera del análisis por diseño. Tras eliminar los casos
con valores perdidos en las variables de los modelos, la muestra
analítica es de 19.731 personas, el 89,4 \% de la muestra elegible. El
apéndice reporta el tamaño muestral por país.

\subsection{Variables}\label{variables}

La variable dependiente son las preferencias frente al cambio
tecnológico, medidas con la batería ``Los gobiernos pueden adoptar
medidas para ayudar a trabajadores e industrias a enfrentar los desafíos
de la digitalización y el cambio tecnológico, como habilidades
obsoletas, escasez de habilidades y posibles pérdidas de empleo.
Teniendo en cuenta cuánto podrían costar y cómo usted y su familia
podrían beneficiarse, ¿en qué medida se opondría o apoyaría que el
gobierno adoptara las siguientes acciones como respuesta a la
digitalización y el cambio tecnológico?''.\footnote{Las formulaciones de
  los ítems corresponden a traducciones propias del cuestionario en
  inglés.} Cada una de las doce medidas se responde en una escala de
cinco puntos, desde muy en contra (1) hasta muy a favor (5), y la opción
``no sabe'' se trata como valor perdido. Aunque la pregunta pide
considerar los costos, cada medida se evalúa por separado y sin una
restricción presupuestaria.

Las medidas se agrupan en tres repertorios. La inversión reúne la
preferencia por la inversión en educación universitaria y formación
profesional para jóvenes, en reentrenamiento para personas en edad de
trabajar y en infraestructura digital. La compensación reúne las
preferencias por prestaciones y servicios más accesibles o generosos,
como el seguro de desempleo o los programas de ingreso mínimo, y una
renta básica universal e incondicional. La protección reúne las
preferencias por la introducción o el aumento de un impuesto a los
robots o a las empresas tecnológicas, los subsidios a las empresas de
las industrias más afectadas por el cambio tecnológico, la prioridad a
los trabajadores del país en la contratación y la introducción de
aranceles para proteger las industrias nacionales. Las tres medidas
restantes, que son la limitación o reducción de la jornada laboral y la
promoción de la migración de trabajadores calificados y no calificados,
no forman parte de los repertorios.

Esta agrupación combina criterios teóricos y empíricos. Un análisis
factorial exploratorio de las doce medidas, con cinco factores según el
análisis paralelo, distingue la inversión, la compensación y la
protección, junto a un factor de migración y otro de reducción de la
jornada. Los subsidios presentan cargas similares en la protección y en
la compensación, por lo que el apéndice reporta también los modelos con
un índice de protección que los excluye. La consistencia interna es
moderada en la inversión (α = 0,66) y en la protección (α = 0,60). En la
compensación, que reúne dos ítems, la correlación entre ítems es de
0,43, con un coeficiente de Spearman-Brown de 0,60
(\citeproc{ref-eisinga_reliability_2013}{\textbf{eisinga\_reliability\_2013?}}).
El análisis factorial completo se presenta en el apéndice.

A partir de estas medidas se construyen dos componentes. El volumen
corresponde al promedio de las doce medidas y expresa el nivel general
de apoyo a la intervención estatal frente al cambio tecnológico. La
dirección se mide como la diferencia entre el promedio de cada
repertorio y el volumen de cada persona, de modo que los valores
positivos indican que el repertorio recibe más apoyo que el conjunto de
las medidas y los negativos, menos. Este procedimiento centra las
respuestas de cada persona en su propio promedio, como propone Schwartz
(\citeproc{ref-schwartz_overview_2012}{2012}) para medir prioridades
valóricas, y separa la prioridad asignada a cada repertorio de la
tendencia general a apoyar las medidas, que puede variar entre personas
y países por razones sustantivas y por estilos de respuesta
(\citeproc{ref-rudnev_caveats_2021}{Rudnev, 2021};
\citeproc{ref-fischer_standardization_2004}{\textbf{fischer\_standardization\_2004?}}).

La Figura~\ref{fig-policies} presenta el apoyo promedio a cada medida
entre las personas ocupadas. Los tres repertorios reciben, en promedio,
un apoyo superior al punto medio de la escala, lo que indica una demanda
extendida de intervención estatal frente al cambio tecnológico. La
inversión concentra el mayor apoyo, en línea con estudios previos
(\citeproc{ref-busemeyer_public_2017}{Busemeyer \& Garritzmann, 2017};
\citeproc{ref-churchill_policy_2025}{Churchill et~al., 2025}).

\begin{figure}

\caption{\label{fig-policies}Apoyo promedio a las medidas frente a la
digitalización y el cambio tecnológico}

\centering{

\centering{

\includegraphics[width=0.85\linewidth,height=\textheight,keepaspectratio]{paper_files/figure-pdf/fig-policies-1.pdf}

}

\subcaption{\label{fig-policies}Fuente: elaboración propia con datos de
la OCDE, encuesta Risks That Matter 2024. Personas ocupadas. Escala de 1
(muy en contra) a 5 (muy a favor)}

}

\end{figure}%

Tres conjuntos de variables independientes constituyen el núcleo del
análisis. El primero corresponde a las percepciones sobre el propio
empleo, medidas con la batería ``¿Qué tan probable cree que es que
ocurra lo siguiente con su trabajo (u oportunidades laborales) en los
próximos cinco años?'', cuyos ítems se responden en una escala de
probabilidad de cuatro puntos. Cuatro ítems se refieren a la amenaza
tecnológica y preguntan si el trabajo será reemplazado por un robot, por
una herramienta de inteligencia artificial como ChatGPT o por una
persona que ofrece un servicio similar a través de una plataforma de
internet, y si la persona perderá su empleo por no manejar bien las
nuevas tecnologías o por ser reemplazada por alguien con mejores
habilidades tecnológicas. Tres ítems se refieren al beneficio esperado y
preguntan si la tecnología hará que el trabajo y la jornada sean más
compatibles con la vida privada, que el trabajo sea menos peligroso o
físicamente exigente y que sea menos aburrido, repetitivo, estresante o
mentalmente exigente. Un análisis factorial exploratorio de los siete
ítems respalda una estructura de dos factores, a partir de la cual se
construyen dos índices como el promedio de sus ítems, la amenaza
tecnológica (α = 0,87) y el beneficio esperado (α = 0,78). El ítem sobre
la inteligencia artificial carga en el factor de amenaza en la solución
de dos factores y aparece como una dimensión aislada en una solución de
tres factores, lo que se interpreta como sobreextracción, por lo que se
mantiene en el índice.

El segundo conjunto corresponde a las expectativas sobre los efectos
sociales de la inteligencia artificial, medidas con la batería
``Pensando en los efectos de la inteligencia artificial (IA) en el
mercado laboral, ¿qué tan probable cree que es lo siguiente en
{[}país{]} en los próximos 5 a 10 años?'', con la misma escala de cuatro
puntos. Un análisis factorial exploratorio de los siete ítems identifica
dos factores. Las expectativas negativas promedian los ítems sobre el
aumento de la desigualdad de ingresos, del desempleo y de la vigilancia
en el trabajo (α = 0,62), y las expectativas positivas, los ítems sobre
más tiempo libre para la mayoría, trabajos menos aburridos, repetitivos,
estresantes o mentalmente exigentes y la creación de más empleos (α =
0,69). El ítem sobre la necesidad de que muchas personas se reentrenen
carga en el factor negativo, pero se excluye porque describe un
requisito de adaptación y no un daño.

El tercer conjunto se refiere al contexto institucional y comprende dos
indicadores nacionales, uno para la compensación y otro para la
inversión. Las prestaciones por desempleo se miden con la proporción de
personas desempleadas que reciben una prestación monetaria por
desempleo. Es el indicador 1.3.1 de los Objetivos de Desarrollo
Sostenible y proviene de la base de datos de protección social de la
Organización Internacional del Trabajo. En varios países alcanza su
máximo de 100 \%, lo que reduce su variación en la parte alta de la
escala. Las políticas de mercado laboral se miden con el gasto público
en políticas activas del mercado laboral como porcentaje del PIB, tomado
de la base de gasto social de la OCDE (SOCX). Este gasto comprende los
servicios públicos de empleo y su administración, la capacitación, la
rotación y el reparto del trabajo, los incentivos a la contratación, el
empleo con apoyo y la rehabilitación, la creación directa de empleo y
los incentivos al emprendimiento. El gasto se divide por la tasa de
desempleo para distinguir el esfuerzo de las políticas de la magnitud
del problema que atienden. Los modelos incluyen además dos controles de
nivel país tomados de los World Development Indicators del Banco
Mundial: el PIB per cápita (PPA, en logaritmo; serie NY.GDP.PCAP.PP.KD),
como control del nivel de desarrollo, y el coeficiente de Gini del
ingreso per cápita de los hogares (serie SI.POV.GINI), como control de
la desigualdad. En los 27 países el Gini proviene de encuestas de
ingreso, con datos de entre 2021 y 2024. Para cada país se usa el valor
más reciente disponible desde 2015, y todos los indicadores se
estandarizan entre los 27 países, de modo que sus coeficientes expresan
la asociación con una diferencia de una desviación estándar entre
países.

Los modelos incluyen como controles la preocupación por perder el empleo
o los ingresos del trabajo independiente en los próximos uno o dos años,
medida en una escala de cuatro puntos desde nada preocupado hasta muy
preocupado, el sexo, el grupo de edad, el nivel educacional y el decil
de ingreso autorreportado. Incluyen también cuatro características del
empleo. La posición en el mercado laboral distingue a los asalariados
con contrato indefinido, a los asalariados con contrato temporal o sin
contrato y a los trabajadores independientes. Los otros tres controles
corresponden al trabajo mediado por plataformas digitales en el empleo
principal, al empleo en el sector público y al sector de actividad.

La Tabla~\ref{tbl-summary} resume estas variables y sus principales
estadísticos descriptivos en la muestra analítica.

\begin{longtable}[]{@{}
  >{\raggedright\arraybackslash}p{(\linewidth - 4\tabcolsep) * \real{0.3136}}
  >{\raggedright\arraybackslash}p{(\linewidth - 4\tabcolsep) * \real{0.4661}}
  >{\raggedright\arraybackslash}p{(\linewidth - 4\tabcolsep) * \real{0.2203}}@{}}
\caption{Principales estadísticos
descriptivos}\label{tbl-summary}\tabularnewline
\toprule\noalign{}
\begin{minipage}[b]{\linewidth}\raggedright
Etiqueta
\end{minipage} & \begin{minipage}[b]{\linewidth}\raggedright
Estadísticas / Valores
\end{minipage} & \begin{minipage}[b]{\linewidth}\raggedright
Frec. (\% sobre válidos)
\end{minipage} \\
\midrule\noalign{}
\endfirsthead
\toprule\noalign{}
\begin{minipage}[b]{\linewidth}\raggedright
Etiqueta
\end{minipage} & \begin{minipage}[b]{\linewidth}\raggedright
Estadísticas / Valores
\end{minipage} & \begin{minipage}[b]{\linewidth}\raggedright
Frec. (\% sobre válidos)
\end{minipage} \\
\midrule\noalign{}
\endhead
\bottomrule\noalign{}
\endlastfoot
Volumen (apoyo general) & \begin{minipage}[t]{\linewidth}\raggedright
Media (d-s) : 3,6 (0,5)\\
min \textless{} mediana \textless{} max:\\
1 \textless{} 3,6 \textless{} 5\strut
\end{minipage} & 155 valores distintos \\
Dirección: inversión & \begin{minipage}[t]{\linewidth}\raggedright
Media (d-s) : 0,3 (0,5)\\
min \textless{} mediana \textless{} max:\\
-2,3 \textless{} 0,3 \textless{} 3,2\strut
\end{minipage} & 388 valores distintos \\
Dirección: compensación & \begin{minipage}[t]{\linewidth}\raggedright
Media (d-s) : 0 (0,7)\\
min \textless{} mediana \textless{} max:\\
-3 \textless{} 0 \textless{} 2,7\strut
\end{minipage} & 292 valores distintos \\
Dirección: protección & \begin{minipage}[t]{\linewidth}\raggedright
Media (d-s) : 0,1 (0,5)\\
min \textless{} mediana \textless{} max:\\
-2,9 \textless{} 0 \textless{} 2,6\strut
\end{minipage} & 462 valores distintos \\
Amenaza tecnológica & \begin{minipage}[t]{\linewidth}\raggedright
Media (d-s) : 2,1 (0,8)\\
min \textless{} mediana \textless{} max:\\
1 \textless{} 2 \textless{} 4\strut
\end{minipage} & 19 valores distintos \\
Beneficio esperado & \begin{minipage}[t]{\linewidth}\raggedright
Media (d-s) : 2,6 (0,8)\\
min \textless{} mediana \textless{} max:\\
1 \textless{} 2,7 \textless{} 4\strut
\end{minipage} & 13 valores distintos \\
Expectativas negativas sobre la IA &
\begin{minipage}[t]{\linewidth}\raggedright
Media (d-s) : 3 (0,6)\\
min \textless{} mediana \textless{} max:\\
1 \textless{} 3 \textless{} 4\strut
\end{minipage} & 13 valores distintos \\
Expectativas positivas sobre la IA &
\begin{minipage}[t]{\linewidth}\raggedright
Media (d-s) : 2,7 (0,7)\\
min \textless{} mediana \textless{} max:\\
1 \textless{} 2,7 \textless{} 4\strut
\end{minipage} & 13 valores distintos \\
Preocupación por perder el empleo &
\begin{minipage}[t]{\linewidth}\raggedright
1. Nada preocupado\\
2. No muy preocupado\\
3. Algo preocupado\\
4. Muy preocupado\strut
\end{minipage} & \begin{minipage}[t]{\linewidth}\raggedright
2971 (15,1\%)\\
5523 (28,0\%)\\
6109 (31,0\%)\\
5128 (26,0\%)\strut
\end{minipage} \\
Sexo & \begin{minipage}[t]{\linewidth}\raggedright
1. Hombre\\
2. Mujer\strut
\end{minipage} & \begin{minipage}[t]{\linewidth}\raggedright
10567 (53,6\%)\\
9164 (46,4\%)\strut
\end{minipage} \\
Grupo de edad & \begin{minipage}[t]{\linewidth}\raggedright
1. 18-24\\
2. 25-34\\
3. 35-44\\
4. 45-54\\
5. 55-64\strut
\end{minipage} & \begin{minipage}[t]{\linewidth}\raggedright
2135 (10,8\%)\\
4642 (23,5\%)\\
4838 (24,5\%)\\
4612 (23,4\%)\\
3504 (17,8\%)\strut
\end{minipage} \\
Nivel educacional & \begin{minipage}[t]{\linewidth}\raggedright
1. Bajo\\
2. Medio\\
3. Alto\strut
\end{minipage} & \begin{minipage}[t]{\linewidth}\raggedright
1523 ( 7,7\%)\\
8089 (41,0\%)\\
10119 (51,3\%)\strut
\end{minipage} \\
Decil de ingreso & \begin{minipage}[t]{\linewidth}\raggedright
1. {[}1{]} First decile\\
2. {[}2{]} Second decile\\
3. {[}3{]} Third decile\\
4. {[}4{]} Fourth decile\\
5. {[}5{]} Fifth decile\\
6. {[}6{]} Sixth decile\\
7. {[}7{]} Seventh decile\\
8. {[}8{]} Eighth decile\\
9. {[}9{]} Ninth decile\\
10. {[}10{]} Tenth decile\strut
\end{minipage} & \begin{minipage}[t]{\linewidth}\raggedright
3309 (16,8\%)\\
1808 ( 9,2\%)\\
1803 ( 9,1\%)\\
2673 (13,5\%)\\
2345 (11,9\%)\\
1836 ( 9,3\%)\\
1530 ( 7,8\%)\\
1448 ( 7,3\%)\\
1317 ( 6,7\%)\\
1662 ( 8,4\%)\strut
\end{minipage} \\
Posición en el mercado laboral &
\begin{minipage}[t]{\linewidth}\raggedright
1. Asalariado con contrato indefinido\\
2. Asalariado con contrato temporal o sin contrato\\
3. Independiente\strut
\end{minipage} & \begin{minipage}[t]{\linewidth}\raggedright
15150 (76,8\%)\\
2471 (12,5\%)\\
2110 (10,7\%)\strut
\end{minipage} \\
Trabajo en plataformas & \begin{minipage}[t]{\linewidth}\raggedright
1. No\\
2. Sí\strut
\end{minipage} & \begin{minipage}[t]{\linewidth}\raggedright
17889 (90,7\%)\\
1842 ( 9,3\%)\strut
\end{minipage} \\
Sector público & \begin{minipage}[t]{\linewidth}\raggedright
1. No\\
2. Sí\strut
\end{minipage} & \begin{minipage}[t]{\linewidth}\raggedright
13463 (68,2\%)\\
6268 (31,8\%)\strut
\end{minipage} \\
Sector de actividad & \begin{minipage}[t]{\linewidth}\raggedright
1. Agricultura y minería\\
2. Manufactura y servicios básicos\\
3. Construcción\\
4. Comercio, transporte y hotelería\\
5. TIC, finanzas y servicios a empresas\\
6. Administración pública, educación y salud\\
7. Otros servicios\\
8. Otro o no declarado\strut
\end{minipage} & \begin{minipage}[t]{\linewidth}\raggedright
445 ( 2,3\%)\\
2264 (11,5\%)\\
1167 ( 5,9\%)\\
2831 (14,3\%)\\
4598 (23,3\%)\\
4474 (22,7\%)\\
2878 (14,6\%)\\
1074 ( 5,4\%)\strut
\end{minipage} \\
\end{longtable}

\subsection{Métodos}\label{muxe9todos}

Los análisis se realizaron en R, versión 4.6.1. Para examinar la
relación entre las percepciones sobre el cambio tecnológico y las
preferencias se estiman modelos lineales multinivel con personas
anidadas en países, mediante la función lmer del paquete lme4
(\citeproc{ref-bates_fitting_2015}{\textbf{bates\_fitting\_2015?}}).
Este enfoque considera la estructura jerárquica de los datos comparados
y descompone la varianza de cada resultado en un componente individual y
uno nacional
(\citeproc{ref-snijders_multilevel_2012}{\textbf{snijders\_multilevel\_2012?}};
\citeproc{ref-hox_multilevel_2017}{\textbf{hox\_multilevel\_2017?}}),
con interceptos aleatorios por país. Se estima un conjunto de modelos
para el volumen y uno para la dirección de cada repertorio, y en cada
caso un modelo nulo establece primero la proporción de la varianza que
se ubica entre países. Todos los modelos incorporan los ponderadores de
postestratificación de la OCDE en el nivel individual.

Para separar las asociaciones dentro y entre países, las cuatro
percepciones se centran en la media de cada país, de modo que sus
coeficientes expresan diferencias entre personas de un mismo país
(\citeproc{ref-hox_multilevel_2017}{\textbf{hox\_multilevel\_2017?}}).
Los modelos se estiman de forma secuencial. El modelo 1 incluye las
cuatro percepciones, el modelo 2 agrega los controles sociodemográficos
y laborales, el modelo 3 la preocupación por perder el empleo y el
modelo 4 los cuatro indicadores de nivel país: el PIB per cápita, el
Gini, las prestaciones por desempleo y las políticas de mercado laboral.
El contexto institucional se incorpora en dos pasos para la inversión y
la compensación. Primero, el modelo 4 incluye las prestaciones por
desempleo, que corresponden a la compensación, y las políticas de
mercado laboral, que corresponden a la inversión. Luego, se estiman las
interacciones entre el indicador de cada repertorio y la amenaza
tecnológica, y entre ese indicador y las expectativas negativas, una
interacción por modelo y con una pendiente aleatoria para la percepción
correspondiente
(\citeproc{ref-heisig_why_2019}{\textbf{heisig\_why\_2019?}}). Todos los
modelos que incorporan el contexto institucional incluyen los cuatro
indicadores de nivel país.

Para examinar si la relación entre la amenaza tecnológica y las
preferencias varía entre grupos, se estiman además modelos con una
interacción entre la amenaza y el grupo de edad, y entre la amenaza y el
nivel educacional. El apéndice presenta, como análisis exploratorio, las
interacciones con el sector de actividad.

Las hipótesis sobre el volumen (H1 y H4) se contrastan en los modelos
del volumen, y las hipótesis sobre la dirección (H2a, H2b, H3, H5a y
H5b), en los modelos de cada repertorio. La hipótesis termostática (H6b)
se evalúa en la Tabla~\ref{tbl-modelos}, con el efecto de cada indicador
sobre su repertorio, y la hipótesis de disponibilidad (H6a), con las
interacciones entre niveles de la Tabla~\ref{tbl-interacciones}. El
apéndice descompone, además, la protección en sus cuatro medidas, para
distinguir las dirigidas a la tecnología de las dirigidas a la
competencia extranjera.

\section{Resultados}\label{resultados}

{

\begin{longtable}[]{@{}
  >{\raggedright\arraybackslash}p{(\linewidth - 8\tabcolsep) * \real{0.3684}}
  >{\centering\arraybackslash}p{(\linewidth - 8\tabcolsep) * \real{0.1579}}
  >{\centering\arraybackslash}p{(\linewidth - 8\tabcolsep) * \real{0.1579}}
  >{\centering\arraybackslash}p{(\linewidth - 8\tabcolsep) * \real{0.1579}}
  >{\centering\arraybackslash}p{(\linewidth - 8\tabcolsep) * \real{0.1579}}@{}}

\caption{\label{tbl-modelos}Modelos multinivel del volumen y la
dirección de las preferencias frente al cambio tecnológico}

\tabularnewline

\toprule\noalign{}
\begin{minipage}[b]{\linewidth}\raggedright
\end{minipage} & \begin{minipage}[b]{\linewidth}\centering
Volumen
\end{minipage} & \begin{minipage}[b]{\linewidth}\centering
Inversión
\end{minipage} & \begin{minipage}[b]{\linewidth}\centering
Compensación
\end{minipage} & \begin{minipage}[b]{\linewidth}\centering
Protección
\end{minipage} \\
\midrule\noalign{}
\endhead
\bottomrule\noalign{}
\endlastfoot
Intercepto & 3,301 *** & 0,161 *** & -0,152 *** & 0,082 ** \\
& (0,035) & (0,035) & (0,046) & (0,031) \\
\textbf{Nivel individual} & & & & \\
Amenaza tecnológica & 0,005 & -0,157 *** & -0,042 *** & 0,009 \\
& (0,005) & (0,005) & (0,007) & (0,005) \\
Beneficio esperado & 0,091 *** & 0,058 *** & 0,024 ** & -0,044 *** \\
& (0,006) & (0,005) & (0,008) & (0,005) \\
Expectativas negativas sobre la IA & 0,192 *** & -0,011 & 0,058 *** &
0,108 *** \\
& (0,006) & (0,006) & (0,008) & (0,005) \\
Expectativas positivas sobre la IA & 0,142 *** & -0,004 & -0,045 *** &
-0,082 *** \\
& (0,006) & (0,006) & (0,008) & (0,006) \\
Preocupación por perder el empleo & 0,042 *** & -0,017 *** & 0,050 *** &
0,007 * \\
& (0,004) & (0,004) & (0,005) & (0,003) \\
Mujer & 0,025 *** & -0,018 * & 0,039 *** & 0,020 ** \\
& (0,007) & (0,007) & (0,010) & (0,007) \\
25 a 34 años & 0,063 *** & 0,037 ** & 0,035 & 0,005 \\
& (0,014) & (0,013) & (0,019) & (0,012) \\
35 a 44 años & 0,065 *** & 0,086 *** & 0,037 * & 0,008 \\
& (0,013) & (0,013) & (0,018) & (0,012) \\
45 a 54 años & 0,065 *** & 0,140 *** & 0,036 & 0,028 * \\
& (0,014) & (0,013) & (0,019) & (0,012) \\
55 a 64 años & 0,071 *** & 0,172 *** & 0,041 * & 0,015 \\
& (0,014) & (0,014) & (0,020) & (0,013) \\
Educación media & -0,009 & 0,069 *** & -0,014 & -0,009 \\
& (0,015) & (0,014) & (0,020) & (0,013) \\
Educación alta & 0,009 & 0,133 *** & -0,071 *** & -0,101 *** \\
& (0,015) & (0,014) & (0,021) & (0,014) \\
Decil de ingreso & -0,002 & 0,008 *** & -0,016 *** & -0,001 \\
& (0,001) & (0,001) & (0,002) & (0,001) \\
Asalariado temporal o sin contrato & 0,012 & 0,028 ** & 0,001 & -0,019
* \\
& (0,011) & (0,010) & (0,015) & (0,010) \\
Independiente & -0,066 *** & 0,013 & 0,001 & -0,028 ** \\
& (0,011) & (0,011) & (0,015) & (0,010) \\
Trabajo en plataformas & 0,122 *** & -0,101 *** & -0,046 ** & 0,012 \\
& (0,013) & (0,012) & (0,017) & (0,011) \\
Sector público & -0,026 ** & -0,048 *** & 0,004 & 0,009 \\
& (0,009) & (0,009) & (0,012) & (0,008) \\
\textbf{Nivel país} & & & & \\
PIB per cápita (log) & 0,005 & -0,003 & 0,008 & -0,010 \\
& (0,022) & (0,025) & (0,023) & (0,017) \\
Gini & 0,037 & 0,029 & -0,011 & 0,004 \\
& (0,022) & (0,024) & (0,022) & (0,017) \\
Prestaciones por desempleo & -0,024 & 0,010 & -0,044 * & 0,001 \\
& (0,021) & (0,024) & (0,022) & (0,016) \\
Políticas de mercado laboral & -0,035 & 0,025 & -0,038 & -0,014 \\
& (0,022) & (0,025) & (0,023) & (0,017) \\
Log verosimilitud & -15080,342 & -14147,586 & -21267,176 & -13104,640 \\
N & 19.731 & 19.731 & 19.731 & 19.731 \\
Países & 27 & 27 & 27 & 27 \\
Varianza del intercepto (país) & 0,007 & 0,009 & 0,007 & 0,004 \\
Varianza residual & 0,239 & 0,217 & 0,448 & 0,196 \\

\end{longtable}

}

\emph{Nota}: Categorías de referencia: hombre, 18 a 24 años, educación
baja, asalariado con contrato indefinido, sin trabajo en plataformas y
no empleado en el sector público (asalariados del sector privado e
independientes). Todos los modelos controlan además por sector de
actividad. Errores estándar entre paréntesis. *** p \textless{} 0,001;
** p \textless{} 0,01; * p \textless{} 0,05. Los indicadores de nivel
país están estandarizados entre los 27 países. PIB per cápita en paridad
de poder adquisitivo, en logaritmo, y coeficiente de Gini, ambos del
Banco Mundial. Prestaciones por desempleo: porcentaje de personas
desempleadas que reciben prestaciones monetarias por desempleo (OIT,
indicador 1.3.1 de los ODS). Políticas de mercado laboral: gasto público
en políticas activas del mercado laboral como porcentaje del PIB,
dividido por la tasa de desempleo (OCDE, SOCX).

{

\begin{longtable}[]{@{}
  >{\raggedright\arraybackslash}p{(\linewidth - 8\tabcolsep) * \real{0.3548}}
  >{\centering\arraybackslash}p{(\linewidth - 8\tabcolsep) * \real{0.1355}}
  >{\centering\arraybackslash}p{(\linewidth - 8\tabcolsep) * \real{0.1677}}
  >{\centering\arraybackslash}p{(\linewidth - 8\tabcolsep) * \real{0.1548}}
  >{\centering\arraybackslash}p{(\linewidth - 8\tabcolsep) * \real{0.1871}}@{}}

\caption{\label{tbl-interacciones}Interacciones entre las percepciones y
los indicadores institucionales}

\tabularnewline

\toprule\noalign{}
\begin{minipage}[b]{\linewidth}\raggedright
\end{minipage} & \begin{minipage}[b]{\linewidth}\centering
Inversión × amenaza
\end{minipage} & \begin{minipage}[b]{\linewidth}\centering
Inversión × expectativas
\end{minipage} & \begin{minipage}[b]{\linewidth}\centering
Compensación × amenaza
\end{minipage} & \begin{minipage}[b]{\linewidth}\centering
Compensación × expectativas
\end{minipage} \\
\midrule\noalign{}
\endhead
\bottomrule\noalign{}
\endlastfoot
Intercepto & 0,155 *** & 0,163 *** & -0,147 ** & -0,152 *** \\
& (0,035) & (0,035) & (0,046) & (0,046) \\
\textbf{Nivel individual} & & & & \\
Amenaza tecnológica & -0,157 *** & -0,158 *** & -0,041 *** & -0,041
*** \\
& (0,009) & (0,005) & (0,012) & (0,007) \\
Beneficio esperado & 0,061 *** & 0,058 *** & 0,022 ** & 0,024 ** \\
& (0,005) & (0,005) & (0,008) & (0,008) \\
Expectativas negativas sobre la IA & -0,010 & -0,010 & 0,058 *** & 0,059
*** \\
& (0,006) & (0,009) & (0,008) & (0,012) \\
Expectativas positivas sobre la IA & -0,004 & -0,005 & -0,045 *** &
-0,044 *** \\
& (0,006) & (0,006) & (0,008) & (0,008) \\
Preocupación por perder el empleo & -0,015 *** & -0,017 *** & 0,049 ***
& 0,050 *** \\
& (0,004) & (0,004) & (0,005) & (0,005) \\
Mujer & -0,020 ** & -0,018 ** & 0,038 *** & 0,039 *** \\
& (0,007) & (0,007) & (0,010) & (0,010) \\
25 a 34 años & 0,033 * & 0,036 ** & 0,038 * & 0,036 \\
& (0,013) & (0,013) & (0,019) & (0,019) \\
35 a 44 años & 0,082 *** & 0,083 *** & 0,039 * & 0,038 * \\
& (0,013) & (0,013) & (0,018) & (0,018) \\
45 a 54 años & 0,134 *** & 0,138 *** & 0,039 * & 0,038 * \\
& (0,013) & (0,013) & (0,019) & (0,019) \\
55 a 64 años & 0,167 *** & 0,170 *** & 0,044 * & 0,042 * \\
& (0,014) & (0,014) & (0,020) & (0,020) \\
Educación media & 0,072 *** & 0,069 *** & -0,017 & -0,013 \\
& (0,014) & (0,014) & (0,020) & (0,020) \\
Educación alta & 0,136 *** & 0,133 *** & -0,076 *** & -0,071 *** \\
& (0,014) & (0,014) & (0,020) & (0,021) \\
Decil de ingreso & 0,008 *** & 0,008 *** & -0,016 *** & -0,016 *** \\
& (0,001) & (0,001) & (0,002) & (0,002) \\
Asalariado temporal o sin contrato & 0,026 * & 0,028 ** & 0,002 &
0,002 \\
& (0,010) & (0,010) & (0,015) & (0,015) \\
Independiente & 0,012 & 0,014 & 0,002 & 0,001 \\
& (0,011) & (0,011) & (0,015) & (0,015) \\
Trabajo en plataformas & -0,101 *** & -0,100 *** & -0,042 * & -0,044
* \\
& (0,012) & (0,012) & (0,017) & (0,017) \\
Sector público & -0,046 *** & -0,047 *** & 0,001 & 0,004 \\
& (0,008) & (0,009) & (0,012) & (0,012) \\
\textbf{Nivel país} & & & & \\
PIB per cápita (log) & -0,016 & -0,007 & 0,003 & 0,007 \\
& (0,023) & (0,023) & (0,018) & (0,023) \\
Gini & 0,025 & 0,021 & -0,002 & -0,014 \\
& (0,023) & (0,023) & (0,017) & (0,022) \\
Prestaciones por desempleo & -0,003 & 0,000 & -0,041 * & -0,044 * \\
& (0,022) & (0,022) & (0,019) & (0,021) \\
Políticas de mercado laboral & 0,038 & 0,027 & -0,027 & -0,038 \\
& (0,024) & (0,024) & (0,017) & (0,022) \\
\textbf{Interacciones entre niveles} & & & & \\
Amenaza tecnológica × Políticas de mercado laboral & -0,031 *** & & & \\
& (0,009) & & & \\
Expectativas negativas × Políticas de mercado laboral & & -0,021 * &
& \\
& & (0,009) & & \\
Amenaza tecnológica × Prestaciones por desempleo & & & 0,047 *** & \\
& & & (0,011) & \\
Expectativas negativas × Prestaciones por desempleo & & & & 0,013 \\
& & & & (0,012) \\
Log verosimilitud & -14101,113 & -14132,162 & -21221,593 & -21263,064 \\
N & 19.731 & 19.731 & 19.731 & 19.731 \\
Países & 27 & 27 & 27 & 27 \\
Varianza del intercepto (país) & 0,009 & 0,008 & 0,006 & 0,007 \\
Varianza de la pendiente de la amenaza tecnológica & 0,002 & & 0,002
& \\
Covarianza entre intercepto y pendiente (amenaza) & -0,002 & & -0,003
& \\
Varianza residual & 0,216 & 0,217 & 0,446 & 0,447 \\
Varianza de la pendiente de las expectativas negativas & & 0,001 & &
0,002 \\
Covarianza entre intercepto y pendiente (expectativas) & & -0,002 & &
-0,000 \\

\end{longtable}

}

\emph{Nota}: Categorías de referencia: hombre, 18 a 24 años, educación
baja, asalariado con contrato indefinido, sin trabajo en plataformas y
no empleado en el sector público (asalariados del sector privado e
independientes). Todos los modelos controlan además por sector de
actividad. Errores estándar entre paréntesis. *** p \textless{} 0,001;
** p \textless{} 0,01; * p \textless{} 0,05. Los indicadores de nivel
país están estandarizados entre los 27 países. PIB per cápita en paridad
de poder adquisitivo, en logaritmo, y coeficiente de Gini, ambos del
Banco Mundial. Prestaciones por desempleo: porcentaje de personas
desempleadas que reciben prestaciones monetarias por desempleo (OIT,
indicador 1.3.1 de los ODS). Políticas de mercado laboral: gasto público
en políticas activas del mercado laboral como porcentaje del PIB,
dividido por la tasa de desempleo (OCDE, SOCX). Cada modelo incluye una
pendiente aleatoria por país para la percepción que interactúa.

\begin{figure}

\caption{\label{fig-interacciones}Asociación de la amenaza tecnológica y
de las expectativas negativas sobre la IA con la dirección de las
preferencias, según el indicador institucional de cada repertorio}

\centering{

\includegraphics[width=0.85\linewidth,height=\textheight,keepaspectratio]{paper_files/figure-pdf/fig-interacciones-1.pdf}

}

\end{figure}%

\begin{figure}[H]

\caption{\label{fig-moderacion}Pendiente de la amenaza tecnológica según
grupo de edad y nivel educacional}

\centering{

\includegraphics[width=0.85\linewidth,height=\textheight,keepaspectratio]{paper_files/figure-pdf/fig-moderacion-1.pdf}

}

\end{figure}%

\emph{Nota}: Pendientes estimadas a partir de modelos con una
interacción entre la amenaza tecnológica y cada moderador. Intervalos de
confianza de 95 \%.

\section{Discusión}\label{discusiuxf3n}

\section{Conclusión}\label{conclusiuxf3n}

\section{Referencias}\label{referencias}

\protect\phantomsection\label{refs}
\begin{CSLReferences}{1}{0}
\bibitem[\citeproctext]{ref-acemoglu_tasks_2022}
Acemoglu, D., \& Restrepo, P. (2022). Tasks, {Automation}, and the
{Rise} in {U}.{S}. {Wage Inequality}. \emph{Econometrica : Journal of
the Econometric Society}, \emph{90}(5), 1973-2016.
\url{https://doi.org/10.3982/ECTA19815}

\bibitem[\citeproctext]{ref-anderson_workers_2007}
Anderson, C. J., \& Pontusson, J. (2007). Workers, Worries and Welfare
States: {Social} Protection and Job Insecurity in 15 {OECD} Countries.
\emph{European Journal of Political Research}, \emph{46}(2), 211-235.
\url{https://doi.org/10.1111/j.1475-6765.2007.00692.x}

\bibitem[\citeproctext]{ref-anelli_individual_2021}
Anelli, M., Colantone, I., \& Stanig, P. (2021). Individual
Vulnerability to Industrial Robot Adoption Increases Support for the
Radical Right. \emph{Proceedings of the National Academy of Sciences},
\emph{118}(47), e2111611118.
\url{https://doi.org/10.1073/pnas.2111611118}

\bibitem[\citeproctext]{ref-arts_three_2002}
Arts, W., \& Gelissen, J. (2002). Three Worlds of Welfare Capitalism or
More? {A} State-of-the-Art Report. \emph{Journal of European Social
Policy}, \emph{12}(2), 137-158.
\url{https://doi.org/10.1177/0952872002012002114}

\bibitem[\citeproctext]{ref-autor_why_2015}
Autor, D. H. (2015). Why {Are There Still So Many Jobs}? {The History}
and {Future} of {Workplace Automation}. \emph{Journal of Economic
Perspectives}, \emph{29}(3), 3-30.
\url{https://doi.org/10.1257/jep.29.3.3}

\bibitem[\citeproctext]{ref-behrendt_social_2019}
Behrendt, C., Nguyen, Q. A., \& Rani, U. (2019). Social Protection
Systems and the Future of Work: {Ensuring} Social Security for Digital
Platform Workers. \emph{International Social Security Review},
\emph{72}(3), 17-41. \url{https://doi.org/10.1111/issr.12212}

\bibitem[\citeproctext]{ref-bicchi_unpacking_2025}
Bicchi, N., Kuo, A., \& Gallego, A. (2025). Unpacking Technological
Risks: {Different} Sources of Concern and Policy Preferences.
\emph{Political Studies}, \emph{73}(3), 1054-1077.
\url{https://doi.org/10.1177/00323217241281575}

\bibitem[\citeproctext]{ref-busemeyer_preferred_2023}
Busemeyer, M. R., Gandenberger, M., Knotz, C., \& Tober, T. (2023).
Preferred Policy Responses to Technological Change: {Survey} Evidence
from {OECD} Countries. \emph{Socio-Economic Review}, \emph{21}(1),
593-615. \url{https://doi.org/10.1093/ser/mwac015}

\bibitem[\citeproctext]{ref-busemeyer_public_2017}
Busemeyer, M. R., \& Garritzmann, J. L. (2017). Public Opinion on Policy
and Budgetary Trade-Offs in {European} Welfare States: {Evidence} from a
New Comparative Survey. \emph{Journal of European Public Policy},
\emph{24}(6), 871-889.
\url{https://doi.org/10.1080/13501763.2017.1298658}

\bibitem[\citeproctext]{ref-busemeyer_social_2022}
Busemeyer, M. R., \& Sahm, A. H. J. (2022). Social {Investment},
{Redistribution} or {Basic Income}? {Exploring} the {Association Between
Automation Risk} and {Welfare State Attitudes} in {Europe}.
\emph{Journal of Social Policy}, \emph{51}(4), 751-770.
\url{https://doi.org/10.1017/S0047279421000519}

\bibitem[\citeproctext]{ref-busemeyer_dealing_2023}
Busemeyer, M. R., \& Tober, T. (2023). Dealing with Technological
Change: {Social} Policy Preferences and Institutional Context.
\emph{Comparative Political Studies}, \emph{56}(7), 968-999.
\url{https://doi.org/10.1177/00104140221139381}

\bibitem[\citeproctext]{ref-cavaille_fair_2025}
Cavaillé, C. (2025). \emph{Fair {Enough}?: {Support} for
{Redistribution} in the {Age} of {Inequality}} (1.ª ed.). Cambridge
University Press. \url{https://doi.org/10.1017/9781009366038}

\bibitem[\citeproctext]{ref-chung_institutions_2011}
Chung, H., \& Van Oorschot, W. (2011). Institutions versus Market
Forces: {Explaining} the Employment Insecurity of {European} Individuals
during (the Beginning of) the Financial Crisis. \emph{Journal of
European Social Policy}, \emph{21}(4), 287-301.
\url{https://doi.org/10.1177/0958928711412224}

\bibitem[\citeproctext]{ref-churchill_policy_2025}
Churchill, B., Vromen, A., \& Ruppanner, L. (2025). Policy Interventions
for the Future of Work: {Do Australians} Want Government Compensation,
Guidance, or Investment to Manage Technological Change? \emph{Policy
Studies}. \url{https://doi.org/10.1080/01442872.2025.2567875}

\bibitem[\citeproctext]{ref-colantone_trade_2018}
Colantone, I., \& Stanig, P. (2018). The Trade Origins of Economic
Nationalism: {Import} Competition and Voting Behavior in {Western
Europe}. \emph{American Journal of Political Science}, \emph{62}(4),
936-953. \url{https://doi.org/10.1111/ajps.12358}

\bibitem[\citeproctext]{ref-cusack_risks_2006}
Cusack, T., Iversen, T., \& Rehm, P. (2006). Risks at {Work}: {The
Demand} and {Supply Sides} of {Government Redistribution}. \emph{Oxford
Review of Economic Policy}, \emph{22}(3), 365-389.
\url{https://doi.org/10.1093/oxrep/grj022}

\bibitem[\citeproctext]{ref-dekker_fear_2017}
Dekker, F., Salomons, A., \& Waal, J. V. D. (2017). Fear of Robots at
Work: {The} Role of Economic Self-Interest. \emph{Socio-Economic
Review}, \emph{15}(3), 539-562. \url{https://doi.org/10.1093/ser/mwx005}

\bibitem[\citeproctext]{ref-eloundou_gpts_2024}
Eloundou, T., Manning, S., Mishkin, P., \& Rock, D. (2024). {GPTs} Are
{GPTs}: {Labor} Market Impact Potential of {LLMs}. \emph{Science},
\emph{384}(6702), 1306-1308.
\url{https://doi.org/10.1126/science.adj0998}

\bibitem[\citeproctext]{ref-emmenegger_age_2012}
Emmenegger, P., Häusermann, S., Palier, B., \& Seeleib-Kaiser, M.
(2012). \emph{The Age of Dualization : {The} Changing Face of Inequality
in Deindustrializing Societies}. (P. Emmenegger, S. Häusermann, B.
Palier, \& M. Seeleib-Kaiser, Eds.). Oxford University Press.
\url{https://doi.org/10.5167/UZH-65793}

\bibitem[\citeproctext]{ref-esping-andersen_three_1993}
Esping-Andersen, G. (1990). \emph{The Three Worlds of Welfare
Capitalism}. Cambridge: Polity press.

\bibitem[\citeproctext]{ref-estevez-abe_social_2001}
Estevez-Abe, M., Iversen, T., \& Soskice, D. (2001). Social Protection
and the Formation of Skills: A Reinterpretation of the Welfare State. En
P. A. Hall \& D. Soskice (Eds.), \emph{Varieties of Capitalism: {The}
Institutional Foundations of Comparative Advantage} (pp. 145-183).
Oxford: Oxford University Press.

\bibitem[\citeproctext]{ref-fan_how_2024}
Fan, Z., Ning, J., \& He, A. J. (2024). How {Does Automation Risk Shape
Social Policy Preference}? {Employment Insecurity} and {Policy Feedback
Effect} in {China}. \emph{Social Policy and Society}, \emph{23}(3),
750-768. \url{https://doi.org/10.1017/S1474746422000513}

\bibitem[\citeproctext]{ref-finseraas_income_2009}
Finseraas, H. (2009). Income {Inequality} and {Demand} for
{Redistribution}: {A Multilevel Analysis} of {European Public Opinion}.
\emph{Scandinavian Political Studies}, \emph{32}(1), 94-119.
\url{https://doi.org/10.1111/j.1467-9477.2008.00211.x}

\bibitem[\citeproctext]{ref-frey_future_2017}
Frey, C. B., \& Osborne, M. A. (2017). The Future of Employment: {How}
Susceptible Are Jobs to Computerisation? \emph{Technological Forecasting
and Social Change}, \emph{114}, 254-280.
\url{https://doi.org/10.1016/j.techfore.2016.08.019}

\bibitem[\citeproctext]{ref-gallego_technological_2022}
Gallego, A., Kuo, A., Manzano, D., \& Fernández-Albertos, J. (2022).
Technological Risk and Policy Preferences. \emph{Comparative Political
Studies}, \emph{55}(1), 60-92.
\url{https://doi.org/10.1177/00104140211024290}

\bibitem[\citeproctext]{ref-gallego_automation_2022}
Gallego, A., \& Kurer, T. (2022). Automation, {Digitalization}, and
{Artificial Intelligence} in the {Workplace}: {Implications} for
{Political Behavior}. \emph{Annual Review of Political Science},
\emph{25}(1), 463-484.
\url{https://doi.org/10.1146/annurev-polisci-051120-104535}

\bibitem[\citeproctext]{ref-garritzmann_public_2018}
Garritzmann, J. L., Busemeyer, M. R., \& Neimanns, E. (2018). Public
Demand for Social Investment: {New} Supporting Coalitions for Welfare
State Reform in {Western Europe}? \emph{Journal of European Public
Policy}, \emph{25}(6), 844-861.
\url{https://doi.org/10.1080/13501763.2017.1401107}

\bibitem[\citeproctext]{ref-gingrich_preferences_2012}
Gingrich, J., \& Ansell, B. (2012). Preferences in Context: {Micro}
Preferences, Macro Contexts, and the Demand for Social Policy.
\emph{Comparative Political Studies}, \emph{45}(12), 1624-1654.
\url{https://doi.org/10.1177/0010414012463904}

\bibitem[\citeproctext]{ref-gmyrek_generative_2023}
Gmyrek, P., Berg, J., \& Bescond, D. (2023). \emph{Generative {AI} and
Jobs: A Global Analysis of Potential Effects on Job Quantity and
Quality} (\{\{ILO Working Paper\}\} No. 96). Geneva: International
Labour Organization. \url{https://doi.org/10.54394/FHEM8239}

\bibitem[\citeproctext]{ref-haslberger_rage_2025}
Haslberger, M., Gingrich, J., \& Bhatia, J. (2025). Rage against the
Machine? {Generative AI} Exposure, Subjective Risk, and Policy
Preferences. \emph{Journal of European Public Policy}.
\url{https://doi.org/10.1080/13501763.2025.2554903}

\bibitem[\citeproctext]{ref-hausermann_politics_2019}
Häusermann, S., Kurer, T., \& Traber, D. (2019). The Politics of
Trade-Offs: {Studying} the Dynamics of Welfare State Reform with
Conjoint Experiments. \emph{Comparative Political Studies},
\emph{52}(7), 1059-1095. \url{https://doi.org/10.1177/0010414018797943}

\bibitem[\citeproctext]{ref-im_losers_2019}
Im, Z. J., Mayer, N., Palier, B., \& Rovny, J. (2019). The {«Losers of
Automation»}: {A} Reservoir of Votes for the Radical Right?
\emph{Research \& Politics}, \emph{6}(1), 2053168018822395.
\url{https://doi.org/10.1177/2053168018822395}

\bibitem[\citeproctext]{ref-iversen_asset_2001}
Iversen, T., \& Soskice, D. (2001). An {Asset Theory} of {Social Policy
Preferences}. \emph{American Political Science Review}, \emph{95}(4),
875-893. \url{https://doi.org/10.1017/S0003055400400079}

\bibitem[\citeproctext]{ref-jaeger_welfare_2006}
Jæger, M. M. (2006). Welfare {Regimes} and {Attitudes Towards
Redistribution}: {The Regime Hypothesis Revisited}. \emph{European
Sociological Review}, \emph{22}(2), 157-170.
\url{https://doi.org/10.1093/esr/jci049}

\bibitem[\citeproctext]{ref-korpi_paradox_1998}
Korpi, W., \& Palme, J. (1998). The {Paradox} of {Redistribution} and
{Strategies} of {Equality}: {Welfare State Institutions}, {Inequality},
and {Poverty} in the {Western Countries}. \emph{American Sociological
Review}, \emph{63}(5), 661. \url{https://doi.org/10.2307/2657333}

\bibitem[\citeproctext]{ref-kuo_automation_2024}
Kuo, A., Manzano, D., \& Gallego, A. (2024). Automation versus Openness:
{Support} for Policies to Address Job Threats. \emph{Journal of Public
Policy}, \emph{44}(1), 1-23.
\url{https://doi.org/10.1017/S0143814X23000260}

\bibitem[\citeproctext]{ref-kurer_declining_2020}
Kurer, T. (2020). The {Declining Middle}: {Occupational Change}, {Social
Status}, and the {Populist Right}. \emph{Comparative Political Studies},
\emph{53}(10-11), 1798-1835.
\url{https://doi.org/10.1177/0010414020912283}

\bibitem[\citeproctext]{ref-kurer_automation_2022}
Kurer, T., \& Häusermann, S. (2022). Automation Risk, Social Policy
Preferences, and Political Participation. En {M. R. Busemeyer, A.
Kemmerling, P. Marx, \& K. van Kersbergen (Eds.)}, \emph{Digitalization
and the Welfare State}. Oxford: Oxford University Press.
\url{https://doi.org/10.1093/oso/9780192848369.003.0008}

\bibitem[\citeproctext]{ref-larsen_institutional_2008}
Larsen, C. A. (2008). The {Institutional Logic} of {Welfare Attitudes}:
{How Welfare Regimes Influence Public Support}. \emph{Comparative
Political Studies}, \emph{41}(2), 145-168.
\url{https://doi.org/10.1177/0010414006295234}

\bibitem[\citeproctext]{ref-lee_luddite_2024}
Lee, J. (2024). Luddite or Technophile? {Policy} Preferences for
Governing Technology-Driven Economic Change. \emph{Socio-Economic
Review}, \emph{22}(3), 1019-1046.
\url{https://doi.org/10.1093/ser/mwae025}

\bibitem[\citeproctext]{ref-marx_effect_2014}
Marx, P. (2014). The Effect of Job Insecurity and Employability on
Preferences for Redistribution in {Western Europe}. \emph{Journal of
European Social Policy}, \emph{24}(4), 351-366.
\url{https://doi.org/10.1177/0958928714538217}

\bibitem[\citeproctext]{ref-mau_moral_2014}
Mau, S. (2014). \emph{The Moral Economy of Welfare States: {Britain} and
{Germany} Compared} (First issued in paperback). London New York NY:
Routledge \& Francis Group.

\bibitem[\citeproctext]{ref-morel_towards_2012}
Morel, N., Palier, B., \& Palme, J. (Eds.). (2012). \emph{Towards a
Social Investment Welfare State? {Ideas}, Policies and Challenges}.
Bristol: Policy Press.

\bibitem[\citeproctext]{ref-oecd_employment_2023}
OECD. (2023). \emph{{OECD Employment Outlook} 2023: {Artificial}
Intelligence and the Labour Market}. Paris: OECD Publishing.
\url{https://doi.org/10.1787/08785bba-en}

\bibitem[\citeproctext]{ref-oecd_risks_2024}
OECD. (2024). Risks That {Matter Survey}: 2024 {Public Use Microdata
File}. {Available} via {http://oe.cd/rtm}.

\bibitem[\citeproctext]{ref-pierson_when_1993}
Pierson, P. (1993). When Effect Becomes Cause: {Policy} Feedback and
Political Change. \emph{World Politics}, \emph{45}(4), 595-628.
\url{https://doi.org/10.2307/2950710}

\bibitem[\citeproctext]{ref-puga_ideology_2023}
Puga, I., \& Moya, C. (2023). Ideology, Legitimation and Collective
Action: {Evidence} from {Chile} on the Mechanism of Ideological
Inversion. \emph{Social Forces}, \emph{101}(3), 1519-1551.
\url{https://doi.org/10.1093/sf/soac032}

\bibitem[\citeproctext]{ref-rehm_risks_2009}
Rehm, P. (2009). Risks and {Redistribution}: {An Individual-Level
Analysis}. \emph{Comparative Political Studies}, \emph{42}(7), 855-881.
\url{https://doi.org/10.1177/0010414008330595}

\bibitem[\citeproctext]{ref-rodrik_why_1998}
Rodrik, D. (1998). Why Do {More Open Economies Have Bigger Governments}?
\emph{Journal of Political Economy}, \emph{106}(5), 997-1032.
\url{https://doi.org/10.1086/250038}

\bibitem[\citeproctext]{ref-rothstein_just_1998}
Rothstein, B. (1998). \emph{Just {Institutions Matter}: {The Moral} and
{Political Logic} of the {Universal Welfare State}} (1.ª ed.). Cambridge
University Press. \url{https://doi.org/10.1017/CBO9780511598449}

\bibitem[\citeproctext]{ref-rudnev_caveats_2021}
Rudnev, M. (2021). Caveats of Non-Ipsatization of Basic Values: {A}
Review of Issues and a Simulation Study. \emph{Journal of Research in
Personality}, \emph{93}, 104118.
\url{https://doi.org/10.1016/j.jrp.2021.104118}

\bibitem[\citeproctext]{ref-rueda_social_2007}
Rueda, D. (2007). \emph{Social {Democracy Inside Out}: {Partisanship}
and {Labor Market Policy} in {Advanced Industrialized Democracies}} (1.ª
ed.). Oxford University PressOxford.
\url{https://doi.org/10.1093/acprof:oso/9780199216352.001.0001}

\bibitem[\citeproctext]{ref-schwander_who_2013}
Schwander, H., \& Häusermann, S. (2013). Who Is in and Who Is out? {A}
Risk-Based Conceptualization of Insiders and Outsiders. \emph{Journal of
European Social Policy}, \emph{23}(3), 248-269.
\url{https://doi.org/10.1177/0958928713480064}

\bibitem[\citeproctext]{ref-schwartz_overview_2012}
Schwartz, S. H. (2012). An {Overview} of the {Schwartz Theory} of {Basic
Values}. \emph{Online Readings in Psychology and Culture}, \emph{2}(1).
\url{https://doi.org/10.9707/2307-0919.1116}

\bibitem[\citeproctext]{ref-soroka_degrees_2010}
Soroka, S. N., \& Wlezien, C. (2010). \emph{Degrees of Democracy:
{Politics}, Public Opinion, and Policy}. Cambridge: Cambridge University
Press.

\bibitem[\citeproctext]{ref-svallfors_worlds_1997}
Svallfors, S. (1997). Worlds of {Welfare} and {Attitudes} to
{Redistribution}: {A Comparison} of {Eight Western Nations}.
\emph{European Sociological Review}, \emph{13}(3), 283-304.
\url{https://doi.org/10.1093/oxfordjournals.esr.a018219}

\bibitem[\citeproctext]{ref-svallfors_government_2013}
Svallfors, Stefan. (2013). Government Quality, Egalitarianism, and
Attitudes to Taxes and Social Spending: A {European} Comparison.
\emph{European Political Science Review}, \emph{5}(3), 363-380.
\url{https://doi.org/10.1017/S175577391200015X}

\bibitem[\citeproctext]{ref-thewissen_automation_2019}
Thewissen, S., \& Rueda, D. (2019). Automation and the {Welfare State}:
{Technological Change} as a {Determinant} of {Redistribution
Preferences}. \emph{Comparative Political Studies}, \emph{52}(2),
171-208. \url{https://doi.org/10.1177/0010414017740600}

\bibitem[\citeproctext]{ref-thomas_survey_2025}
Thomas, J., Frey, V., \& Clarke, C. (2025). \emph{Survey Design and
Technical Documentation Supporting the {OECD Risks} That {Matter
Survey}} (\{\{OECD Social\}\}, \{\{Employment\}\} and \{\{Migration
Working Papers\}\}, 324.ª ed.).
\url{https://doi.org/10.1787/5eebe551-en}

\bibitem[\citeproctext]{ref-vallas_what_2020}
Vallas, S., \& Schor, J. B. (2020). What {Do Platforms Do}?
{Understanding} the {Gig Economy}. \emph{Annual Review of Sociology},
\emph{46}(1), 273-294.
\url{https://doi.org/10.1146/annurev-soc-121919-054857}

\bibitem[\citeproctext]{ref-vanoorschot_making_2006}
Van Oorschot, W. (2006). Making the Difference in Social {Europe}:
{Deservingness} Perceptions among Citizens of {European} Welfare States.
\emph{Journal of European Social Policy}, \emph{16}(1), 23-42.
\url{https://doi.org/10.1177/0958928706059829}

\bibitem[\citeproctext]{ref-walter_globalization_2010}
Walter, S. (2010). Globalization and the Welfare State: {Testing} the
Microfoundations of the Compensation Hypothesis. \emph{International
Studies Quarterly}, \emph{54}(2), 403-426.
\url{https://doi.org/10.1111/j.1468-2478.2010.00593.x}

\bibitem[\citeproctext]{ref-wlezien_public_1995a}
Wlezien, C. (1995). The {Public} as {Thermostat}: {Dynamics} of
{Preferences} for {Spending}. \emph{American Journal of Political
Science}, \emph{39}(4), 981. \url{https://doi.org/10.2307/2111666}

\bibitem[\citeproctext]{ref-wu_misattributed_2022}
Wu, N. (2022). Misattributed Blame? {Attitudes} toward Globalization in
the Age of Automation. \emph{Political Science Research and Methods},
\emph{10}(3), 470-487. \url{https://doi.org/10.1017/psrm.2021.43}

\end{CSLReferences}

\section{Apéndice}\label{apuxe9ndice}

{

\begin{longtable}[]{@{}
  >{\raggedright\arraybackslash}p{(\linewidth - 10\tabcolsep) * \real{0.1471}}
  >{\raggedleft\arraybackslash}p{(\linewidth - 10\tabcolsep) * \real{0.0392}}
  >{\raggedleft\arraybackslash}p{(\linewidth - 10\tabcolsep) * \real{0.1471}}
  >{\raggedleft\arraybackslash}p{(\linewidth - 10\tabcolsep) * \real{0.1176}}
  >{\raggedleft\arraybackslash}p{(\linewidth - 10\tabcolsep) * \real{0.2647}}
  >{\raggedleft\arraybackslash}p{(\linewidth - 10\tabcolsep) * \real{0.2843}}@{}}

\caption{\label{tbl-a-paises}Tamaño de la muestra analítica e
indicadores nacionales por país}

\tabularnewline

\toprule\noalign{}
\begin{minipage}[b]{\linewidth}\raggedright
País
\end{minipage} & \begin{minipage}[b]{\linewidth}\raggedleft
n
\end{minipage} & \begin{minipage}[b]{\linewidth}\raggedleft
PIB per cápita
\end{minipage} & \begin{minipage}[b]{\linewidth}\raggedleft
Gini
\end{minipage} & \begin{minipage}[b]{\linewidth}\raggedleft
Prestaciones por desempleo
\end{minipage} & \begin{minipage}[b]{\linewidth}\raggedleft
Políticas de mercado laboral
\end{minipage} \\
\midrule\noalign{}
\endhead
\bottomrule\noalign{}
\endlastfoot
Alemania & 765 & 62.655 (2024) & 33,7 (2022) & 100,0 (2022) & 0,1564
(2021) \\
Austria & 788 & 63.788 (2024) & 31,2 (2023) & 99,0 (2021) & 0,1341
(2021) \\
Bélgica & 697 & 63.311 (2024) & 26,8 (2023) & 100,0 (2022) & 0,1386
(2021) \\
Canadá & 791 & 57.534 (2024) & 31,5 (2022) & 47,4 (2022) & 0,0405
(2022) \\
Chile & 656 & 30.268 (2024) & 43,0 (2024) & 24,6 (2022) & 0,0149
(2023) \\
Corea & 719 & 55.071 (2024) & 32,9 (2021) & 100,0 (2021) & 0,1481
(2022) \\
Dinamarca & 700 & 71.035 (2024) & 29,9 (2023) & 95,6 (2022) & 0,3242
(2021) \\
Eslovenia & 806 & 48.658 (2024) & 24,7 (2023) & 39,1 (2021) & 0,0478
(2021) \\
España & 672 & 48.460 (2024) & 33,4 (2023) & 57,2 (2022) & 0,0619
(2021) \\
Estados Unidos & 683 & 75.698 (2024) & 41,8 (2024) & 50,5 (2021) &
0,0214 (2023) \\
Estonia & 777 & 41.186 (2024) & 30,7 (2023) & 44,0 (2022) & 0,0889
(2021) \\
Finlandia & 707 & 55.901 (2024) & 27,4 (2023) & 100,0 (2021) & 0,1100
(2021) \\
Francia & 734 & 54.799 (2024) & 31,8 (2023) & 100,0 (2021) & 0,1143
(2022) \\
Grecia & 718 & 37.474 (2024) & 33,4 (2023) & 22,3 (2023) & 0,0209
(2021) \\
Irlanda & 815 & 118.833 (2024) & 29,0 (2023) & 100,0 (2022) & 0,1914
(2021) \\
Israel & 754 & 47.389 (2024) & 38,3 (2022) & 42,8 (2022) & 0,0389
(2023) \\
Italia & 694 & 53.285 (2024) & 34,3 (2023) & 51,5 (2022) & 0,0465
(2021) \\
Letonia & 726 & 37.615 (2024) & 34,0 (2023) & 46,1 (2022) & 0,0184
(2021) \\
Lituania & 733 & 47.462 (2024) & 36,0 (2023) & 88,2 (2022) & 0,0520
(2021) \\
México & 631 & 21.970 (2024) & 42,6 (2024) & 8,3 (2022) & 0,0004
(2023) \\
Noruega & 776 & 94.804 (2024) & 26,5 (2023) & 43,3 (2022) & 0,0865
(2021) \\
Países Bajos & 816 & 70.494 (2024) & 25,7 (2021) & 56,7 (2022) & 0,3585
(2021) \\
Polonia & 793 & 45.153 (2024) & 28,5 (2023) & 31,2 (2020) & 0,1600
(2021) \\
Portugal & 692 & 42.228 (2024) & 33,9 (2023) & 62,3 (2022) & 0,1024
(2021) \\
Reino Unido & 761 & 53.412 (2024) & 32,4 (2021) & 94,1 (2022) & 0,0418
(2019) \\
Suiza & 766 & 85.448 (2024) & 33,5 (2023) & 55,4 (2021) & 0,1209
(2021) \\
Türkiye & 561 & 36.154 (2024) & 43,7 (2023) & 2,5 (2022) & 0,0012
(2021) \\

\end{longtable}

}

\emph{Nota}: Valores en escala original; entre paréntesis, el año del
dato. PIB per cápita en dólares internacionales constantes de 2021
(PPA). Prestaciones por desempleo: porcentaje de personas desempleadas
que reciben prestaciones monetarias por desempleo. Políticas de mercado
laboral: gasto público en políticas activas del mercado laboral como
porcentaje del PIB, dividido por la tasa de desempleo.

{

\begin{longtable}[]{@{}
  >{\raggedright\arraybackslash}p{(\linewidth - 6\tabcolsep) * \real{0.2959}}
  >{\raggedright\arraybackslash}p{(\linewidth - 6\tabcolsep) * \real{0.3163}}
  >{\raggedright\arraybackslash}p{(\linewidth - 6\tabcolsep) * \real{0.2857}}
  >{\raggedright\arraybackslash}p{(\linewidth - 6\tabcolsep) * \real{0.1020}}@{}}

\caption{\label{tbl-a-fuentes}Fuentes de los indicadores de nivel país}

\tabularnewline

\toprule\noalign{}
\begin{minipage}[b]{\linewidth}\raggedright
Indicador
\end{minipage} & \begin{minipage}[b]{\linewidth}\raggedright
Fuente
\end{minipage} & \begin{minipage}[b]{\linewidth}\raggedright
Serie
\end{minipage} & \begin{minipage}[b]{\linewidth}\raggedright
Años
\end{minipage} \\
\midrule\noalign{}
\endhead
\bottomrule\noalign{}
\endlastfoot
PIB per cápita & Banco Mundial, WDI & NY.GDP.PCAP.PP.KD & 2024 \\
Gini & Banco Mundial, WDI & SI.POV.GINI & 2021-2024 \\
Prestaciones por desempleo & OIT, ILOSTAT & SDG\_0131\_RT,
SOC\_CONTIG\_UNE & 2020-2023 \\
Políticas de mercado laboral & OCDE, SOCX; Banco Mundial, WDI & SOCX
TP60; SL.UEM.TOTL.ZS & 2019-2023 \\

\end{longtable}

}

\emph{Nota}: WDI: World Development Indicators. Las definiciones de los
indicadores están en la nota de la Tabla~\ref{tbl-a-paises}. Años: rango
de los años del dato más reciente desde 2015 entre los 27 países.

{

\begin{longtable}[]{@{}
  >{\raggedright\arraybackslash}p{(\linewidth - 4\tabcolsep) * \real{0.3333}}
  >{\raggedright\arraybackslash}p{(\linewidth - 4\tabcolsep) * \real{0.3333}}
  >{\raggedright\arraybackslash}p{(\linewidth - 4\tabcolsep) * \real{0.3333}}@{}}

\caption{\label{tbl-a-afe-q25}Análisis factorial exploratorio de la
amenaza y el beneficio tecnológico en el propio empleo}

\tabularnewline

\toprule\noalign{}
\endhead
\bottomrule\noalign{}
\endlastfoot
~ & Factor 1 & Factor 2 \\
Reemplazo por un robot & 0.86 & -0.05 \\
Reemplazo por una herramienta de IA & 0.83 & 0.00 \\
\begin{minipage}[t]{\linewidth}\raggedright
Reemplazo por una persona en una\\
plataforma de internet\strut
\end{minipage} & 0.76 & 0.06 \\
\begin{minipage}[t]{\linewidth}\raggedright
Pérdida del empleo por habilidades\\
tecnológicas\strut
\end{minipage} & 0.71 & 0.01 \\
Mayor compatibilidad con la vida privada & -0.04 & 0.76 \\
\begin{minipage}[t]{\linewidth}\raggedright
Trabajo menos peligroso o físicamente\\
exigente\strut
\end{minipage} & 0.13 & 0.62 \\
\begin{minipage}[t]{\linewidth}\raggedright
Trabajo menos aburrido, repetitivo o\\
estresante\strut
\end{minipage} & -0.04 & 0.79 \\
Cronbach\textquotesingle s α & 0.88 & 0.78 \\

\end{longtable}

}

{

\begin{longtable}[]{@{}
  >{\raggedright\arraybackslash}p{(\linewidth - 4\tabcolsep) * \real{0.3333}}
  >{\raggedright\arraybackslash}p{(\linewidth - 4\tabcolsep) * \real{0.3333}}
  >{\raggedright\arraybackslash}p{(\linewidth - 4\tabcolsep) * \real{0.3333}}@{}}

\caption{\label{tbl-a-afe-q26}Análisis factorial exploratorio de las
expectativas sobre los efectos sociales de la IA}

\tabularnewline

\toprule\noalign{}
\endhead
\bottomrule\noalign{}
\endlastfoot
~ & Factor 1 & Factor 2 \\
Aumento de la desigualdad de ingresos & 0.63 & -0.03 \\
Más tiempo libre para la mayoría & 0.11 & 0.61 \\
\begin{minipage}[t]{\linewidth}\raggedright
Trabajos menos aburridos, repetitivos o\\
estresantes\strut
\end{minipage} & 0.03 & 0.69 \\
Creación de más empleos & -0.13 & 0.67 \\
Más desempleo & 0.70 & -0.18 \\
Más vigilancia en el trabajo & 0.49 & 0.09 \\
\begin{minipage}[t]{\linewidth}\raggedright
Necesidad de que muchas personas se\\
reentrenen\strut
\end{minipage} & 0.51 & 0.12 \\
Cronbach\textquotesingle s α & 0.67 & 0.69 \\

\end{longtable}

}

{

\begin{longtable}[]{@{}
  >{\raggedright\arraybackslash}p{(\linewidth - 10\tabcolsep) * \real{0.1667}}
  >{\raggedright\arraybackslash}p{(\linewidth - 10\tabcolsep) * \real{0.1667}}
  >{\raggedright\arraybackslash}p{(\linewidth - 10\tabcolsep) * \real{0.1667}}
  >{\raggedright\arraybackslash}p{(\linewidth - 10\tabcolsep) * \real{0.1667}}
  >{\raggedright\arraybackslash}p{(\linewidth - 10\tabcolsep) * \real{0.1667}}
  >{\raggedright\arraybackslash}p{(\linewidth - 10\tabcolsep) * \real{0.1667}}@{}}

\caption{\label{tbl-a-afe-q27}Análisis factorial exploratorio de las
medidas frente al cambio tecnológico}

\tabularnewline

\toprule\noalign{}
\endhead
\bottomrule\noalign{}
\endlastfoot
~ & Factor 1 & Factor 2 & Factor 3 & Factor 4 & Factor 5 \\
\begin{minipage}[t]{\linewidth}\raggedright
Educación universitaria y formación\\
profesional para jóvenes\strut
\end{minipage} & 0.60 & -0.00 & 0.10 & -0.01 & -0.03 \\
\begin{minipage}[t]{\linewidth}\raggedright
Reentrenamiento para personas en edad de\\
trabajar\strut
\end{minipage} & 0.64 & -0.04 & 0.06 & 0.01 & -0.03 \\
Infraestructura digital & 0.60 & 0.11 & -0.11 & -0.05 & 0.13 \\
\begin{minipage}[t]{\linewidth}\raggedright
Impuesto a robots o empresas\\
tecnológicas\strut
\end{minipage} & -0.08 & 0.12 & 0.12 & 0.39 & 0.07 \\
Límite o reducción de la jornada laboral & 0.09 & -0.00 & 0.16 & 0.04 &
0.40 \\
Prestaciones más accesibles o generosas & 0.09 & -0.06 & 0.69 & -0.06 &
-0.03 \\
Renta básica universal e incondicional & -0.06 & -0.02 & 0.66 & -0.06 &
0.13 \\
\begin{minipage}[t]{\linewidth}\raggedright
Subsidios a industrias afectadas por el\\
cambio tecnológico\strut
\end{minipage} & 0.14 & 0.07 & 0.28 & 0.30 & -0.09 \\
Promoción de la migración calificada & 0.23 & 0.72 & -0.07 & -0.03 &
-0.04 \\
Promoción de la migración no calificada & -0.13 & 0.74 & 0.01 & 0.11 &
0.04 \\
\begin{minipage}[t]{\linewidth}\raggedright
Prioridad a trabajadores del país en la\\
contratación\strut
\end{minipage} & 0.16 & -0.30 & -0.08 & 0.46 & 0.06 \\
\begin{minipage}[t]{\linewidth}\raggedright
Aranceles para proteger industrias\\
nacionales\strut
\end{minipage} & -0.04 & 0.07 & -0.10 & 0.74 & -0.03 \\
Cronbach\textquotesingle s α & 0.66 & 0.68 & 0.60 & 0.60 & \\

\end{longtable}

}

{

\begin{longtable}[]{@{}
  >{\raggedright\arraybackslash}p{(\linewidth - 8\tabcolsep) * \real{0.3723}}
  >{\centering\arraybackslash}p{(\linewidth - 8\tabcolsep) * \real{0.1489}}
  >{\centering\arraybackslash}p{(\linewidth - 8\tabcolsep) * \real{0.1596}}
  >{\centering\arraybackslash}p{(\linewidth - 8\tabcolsep) * \real{0.1596}}
  >{\centering\arraybackslash}p{(\linewidth - 8\tabcolsep) * \real{0.1596}}@{}}

\caption{\label{tbl-a-volumen}Modelos secuenciales del volumen}

\tabularnewline

\toprule\noalign{}
\begin{minipage}[b]{\linewidth}\raggedright
\end{minipage} & \begin{minipage}[b]{\linewidth}\centering
Modelo 1
\end{minipage} & \begin{minipage}[b]{\linewidth}\centering
Modelo 2
\end{minipage} & \begin{minipage}[b]{\linewidth}\centering
Modelo 3
\end{minipage} & \begin{minipage}[b]{\linewidth}\centering
Modelo 4
\end{minipage} \\
\midrule\noalign{}
\endhead
\bottomrule\noalign{}
\endlastfoot
Intercepto & 3,566 *** & 3,406 *** & 3,300 *** & 3,301 *** \\
& (0,022) & (0,037) & (0,037) & (0,035) \\
\textbf{Nivel individual} & & & & \\
Amenaza tecnológica & 0,019 *** & 0,017 ** & 0,005 & 0,005 \\
& (0,005) & (0,005) & (0,005) & (0,005) \\
Beneficio esperado & 0,092 *** & 0,093 *** & 0,091 *** & 0,091 *** \\
& (0,006) & (0,006) & (0,006) & (0,006) \\
Expectativas negativas sobre la IA & 0,200 *** & 0,198 *** & 0,192 *** &
0,192 *** \\
& (0,006) & (0,006) & (0,006) & (0,006) \\
Expectativas positivas sobre la IA & 0,142 *** & 0,139 *** & 0,142 *** &
0,142 *** \\
& (0,006) & (0,006) & (0,006) & (0,006) \\
Preocupación por perder el empleo & & & 0,042 *** & 0,042 *** \\
& & & (0,004) & (0,004) \\
Mujer & & 0,032 *** & 0,025 *** & 0,025 *** \\
& & (0,007) & (0,007) & (0,007) \\
25 a 34 años & & 0,068 *** & 0,064 *** & 0,063 *** \\
& & (0,014) & (0,014) & (0,014) \\
35 a 44 años & & 0,071 *** & 0,065 *** & 0,065 *** \\
& & (0,014) & (0,013) & (0,013) \\
45 a 54 años & & 0,071 *** & 0,065 *** & 0,065 *** \\
& & (0,014) & (0,014) & (0,014) \\
55 a 64 años & & 0,073 *** & 0,071 *** & 0,071 *** \\
& & (0,014) & (0,014) & (0,014) \\
Educación media & & -0,006 & -0,009 & -0,009 \\
& & (0,015) & (0,015) & (0,015) \\
Educación alta & & 0,013 & 0,009 & 0,009 \\
& & (0,015) & (0,015) & (0,015) \\
Decil de ingreso & & -0,004 ** & -0,002 & -0,002 \\
& & (0,001) & (0,001) & (0,001) \\
Asalariado temporal o sin contrato & & 0,019 & 0,012 & 0,012 \\
& & (0,011) & (0,011) & (0,011) \\
Independiente & & -0,063 *** & -0,066 *** & -0,066 *** \\
& & (0,011) & (0,011) & (0,011) \\
Trabajo en plataformas & & 0,126 *** & 0,122 *** & 0,122 *** \\
& & (0,013) & (0,013) & (0,013) \\
Sector público & & -0,032 *** & -0,026 ** & -0,026 ** \\
& & (0,009) & (0,009) & (0,009) \\
\textbf{Nivel país} & & & & \\
PIB per cápita (log) & & & & 0,005 \\
& & & & (0,022) \\
Gini & & & & 0,037 \\
& & & & (0,022) \\
Prestaciones por desempleo & & & & -0,024 \\
& & & & (0,021) \\
Políticas de mercado laboral & & & & -0,035 \\
& & & & (0,022) \\
Log verosimilitud & -15183,987 & -15134,592 & -15076,953 & -15080,342 \\
N & 19.731 & 19.731 & 19.731 & 19.731 \\
Países & 27 & 27 & 27 & 27 \\
Varianza del intercepto (país) & 0,013 & 0,013 & 0,011 & 0,007 \\
Varianza residual & 0,243 & 0,241 & 0,239 & 0,239 \\

\end{longtable}

}

\emph{Nota}: Categorías de referencia: hombre, 18 a 24 años, educación
baja, asalariado con contrato indefinido, sin trabajo en plataformas y
no empleado en el sector público (asalariados del sector privado e
independientes). Los modelos 2 a 4 controlan además por sector de
actividad. Errores estándar entre paréntesis. *** p \textless{} 0,001;
** p \textless{} 0,01; * p \textless{} 0,05.

{

\begin{longtable}[]{@{}
  >{\raggedright\arraybackslash}p{(\linewidth - 8\tabcolsep) * \real{0.3684}}
  >{\centering\arraybackslash}p{(\linewidth - 8\tabcolsep) * \real{0.1579}}
  >{\centering\arraybackslash}p{(\linewidth - 8\tabcolsep) * \real{0.1579}}
  >{\centering\arraybackslash}p{(\linewidth - 8\tabcolsep) * \real{0.1579}}
  >{\centering\arraybackslash}p{(\linewidth - 8\tabcolsep) * \real{0.1579}}@{}}

\caption{\label{tbl-a-inversion}Modelos secuenciales de la dirección:
inversión}

\tabularnewline

\toprule\noalign{}
\begin{minipage}[b]{\linewidth}\raggedright
\end{minipage} & \begin{minipage}[b]{\linewidth}\centering
Modelo 1
\end{minipage} & \begin{minipage}[b]{\linewidth}\centering
Modelo 2
\end{minipage} & \begin{minipage}[b]{\linewidth}\centering
Modelo 3
\end{minipage} & \begin{minipage}[b]{\linewidth}\centering
Modelo 4
\end{minipage} \\
\midrule\noalign{}
\endhead
\bottomrule\noalign{}
\endlastfoot
Intercepto & 0,348 *** & 0,118 *** & 0,161 *** & 0,161 *** \\
& (0,018) & (0,033) & (0,034) & (0,035) \\
\textbf{Nivel individual} & & & & \\
Amenaza tecnológica & -0,184 *** & -0,162 *** & -0,157 *** & -0,157
*** \\
& (0,005) & (0,005) & (0,005) & (0,005) \\
Beneficio esperado & 0,063 *** & 0,057 *** & 0,058 *** & 0,058 *** \\
& (0,006) & (0,005) & (0,005) & (0,005) \\
Expectativas negativas sobre la IA & -0,009 & -0,013 * & -0,011 &
-0,011 \\
& (0,006) & (0,006) & (0,006) & (0,006) \\
Expectativas positivas sobre la IA & -0,007 & -0,003 & -0,004 &
-0,004 \\
& (0,006) & (0,006) & (0,006) & (0,006) \\
Preocupación por perder el empleo & & & -0,017 *** & -0,017 *** \\
& & & (0,004) & (0,004) \\
Mujer & & -0,020 ** & -0,018 * & -0,018 * \\
& & (0,007) & (0,007) & (0,007) \\
25 a 34 años & & 0,035 ** & 0,037 ** & 0,037 ** \\
& & (0,013) & (0,013) & (0,013) \\
35 a 44 años & & 0,083 *** & 0,085 *** & 0,086 *** \\
& & (0,013) & (0,013) & (0,013) \\
45 a 54 años & & 0,137 *** & 0,140 *** & 0,140 *** \\
& & (0,013) & (0,013) & (0,013) \\
55 a 64 años & & 0,171 *** & 0,172 *** & 0,172 *** \\
& & (0,014) & (0,014) & (0,014) \\
Educación media & & 0,068 *** & 0,069 *** & 0,069 *** \\
& & (0,014) & (0,014) & (0,014) \\
Educación alta & & 0,131 *** & 0,133 *** & 0,133 *** \\
& & (0,014) & (0,014) & (0,014) \\
Decil de ingreso & & 0,008 *** & 0,008 *** & 0,008 *** \\
& & (0,001) & (0,001) & (0,001) \\
Asalariado temporal o sin contrato & & 0,025 * & 0,028 ** & 0,028 ** \\
& & (0,010) & (0,010) & (0,010) \\
Independiente & & 0,012 & 0,013 & 0,013 \\
& & (0,011) & (0,011) & (0,011) \\
Trabajo en plataformas & & -0,103 *** & -0,101 *** & -0,101 *** \\
& & (0,012) & (0,012) & (0,012) \\
Sector público & & -0,046 *** & -0,048 *** & -0,048 *** \\
& & (0,009) & (0,009) & (0,009) \\
\textbf{Nivel país} & & & & \\
PIB per cápita (log) & & & & -0,003 \\
& & & & (0,025) \\
Gini & & & & 0,029 \\
& & & & (0,024) \\
Prestaciones por desempleo & & & & 0,010 \\
& & & & (0,024) \\
Políticas de mercado laboral & & & & 0,025 \\
& & & & (0,025) \\
Log verosimilitud & -14467,882 & -14143,647 & -14136,976 & -14147,586 \\
N & 19.731 & 19.731 & 19.731 & 19.731 \\
Países & 27 & 27 & 27 & 27 \\
Varianza del intercepto (país) & 0,008 & 0,008 & 0,008 & 0,009 \\
Varianza residual & 0,226 & 0,218 & 0,217 & 0,217 \\

\end{longtable}

}

\emph{Nota}: Categorías de referencia: hombre, 18 a 24 años, educación
baja, asalariado con contrato indefinido, sin trabajo en plataformas y
no empleado en el sector público (asalariados del sector privado e
independientes). Los modelos 2 a 4 controlan además por sector de
actividad. Errores estándar entre paréntesis. *** p \textless{} 0,001;
** p \textless{} 0,01; * p \textless{} 0,05.

{

\begin{longtable}[]{@{}
  >{\raggedright\arraybackslash}p{(\linewidth - 8\tabcolsep) * \real{0.3684}}
  >{\centering\arraybackslash}p{(\linewidth - 8\tabcolsep) * \real{0.1579}}
  >{\centering\arraybackslash}p{(\linewidth - 8\tabcolsep) * \real{0.1579}}
  >{\centering\arraybackslash}p{(\linewidth - 8\tabcolsep) * \real{0.1579}}
  >{\centering\arraybackslash}p{(\linewidth - 8\tabcolsep) * \real{0.1579}}@{}}

\caption{\label{tbl-a-compensacion}Modelos secuenciales de la dirección:
compensación}

\tabularnewline

\toprule\noalign{}
\begin{minipage}[b]{\linewidth}\raggedright
\end{minipage} & \begin{minipage}[b]{\linewidth}\centering
Modelo 1
\end{minipage} & \begin{minipage}[b]{\linewidth}\centering
Modelo 2
\end{minipage} & \begin{minipage}[b]{\linewidth}\centering
Modelo 3
\end{minipage} & \begin{minipage}[b]{\linewidth}\centering
Modelo 4
\end{minipage} \\
\midrule\noalign{}
\endhead
\bottomrule\noalign{}
\endlastfoot
Intercepto & -0,023 & -0,024 & -0,153 ** & -0,152 *** \\
& (0,021) & (0,046) & (0,047) & (0,046) \\
\textbf{Nivel individual} & & & & \\
Amenaza tecnológica & -0,026 *** & -0,029 *** & -0,042 *** & -0,042
*** \\
& (0,007) & (0,007) & (0,007) & (0,007) \\
Beneficio esperado & 0,009 & 0,027 *** & 0,024 ** & 0,024 ** \\
& (0,008) & (0,008) & (0,008) & (0,008) \\
Expectativas negativas sobre la IA & 0,071 *** & 0,066 *** & 0,058 *** &
0,058 *** \\
& (0,008) & (0,008) & (0,008) & (0,008) \\
Expectativas positivas sobre la IA & -0,055 *** & -0,048 *** & -0,045
*** & -0,045 *** \\
& (0,008) & (0,008) & (0,008) & (0,008) \\
Preocupación por perder el empleo & & & 0,051 *** & 0,050 *** \\
& & & (0,005) & (0,005) \\
Mujer & & 0,047 *** & 0,039 *** & 0,039 *** \\
& & (0,010) & (0,010) & (0,010) \\
25 a 34 años & & 0,041 * & 0,036 & 0,035 \\
& & (0,019) & (0,019) & (0,019) \\
35 a 44 años & & 0,044 * & 0,037 * & 0,037 * \\
& & (0,019) & (0,018) & (0,018) \\
45 a 54 años & & 0,044 * & 0,037 & 0,036 \\
& & (0,019) & (0,019) & (0,019) \\
55 a 64 años & & 0,044 * & 0,041 * & 0,041 * \\
& & (0,020) & (0,020) & (0,020) \\
Educación media & & -0,010 & -0,015 & -0,014 \\
& & (0,020) & (0,020) & (0,020) \\
Educación alta & & -0,067 ** & -0,072 *** & -0,071 *** \\
& & (0,021) & (0,021) & (0,021) \\
Decil de ingreso & & -0,018 *** & -0,016 *** & -0,016 *** \\
& & (0,002) & (0,002) & (0,002) \\
Asalariado temporal o sin contrato & & 0,010 & 0,001 & 0,001 \\
& & (0,015) & (0,015) & (0,015) \\
Independiente & & 0,005 & 0,002 & 0,001 \\
& & (0,015) & (0,015) & (0,015) \\
Trabajo en plataformas & & -0,041 * & -0,045 ** & -0,046 ** \\
& & (0,017) & (0,017) & (0,017) \\
Sector público & & -0,003 & 0,004 & 0,004 \\
& & (0,012) & (0,012) & (0,012) \\
\textbf{Nivel país} & & & & \\
PIB per cápita (log) & & & & 0,008 \\
& & & & (0,023) \\
Gini & & & & -0,011 \\
& & & & (0,022) \\
Prestaciones por desempleo & & & & -0,044 * \\
& & & & (0,022) \\
Políticas de mercado laboral & & & & -0,038 \\
& & & & (0,023) \\
Log verosimilitud & -21399,552 & -21305,664 & -21260,985 & -21267,176 \\
N & 19.731 & 19.731 & 19.731 & 19.731 \\
Países & 27 & 27 & 27 & 27 \\
Varianza del intercepto (país) & 0,011 & 0,011 & 0,009 & 0,007 \\
Varianza residual & 0,457 & 0,450 & 0,448 & 0,448 \\

\end{longtable}

}

\emph{Nota}: Categorías de referencia: hombre, 18 a 24 años, educación
baja, asalariado con contrato indefinido, sin trabajo en plataformas y
no empleado en el sector público (asalariados del sector privado e
independientes). Los modelos 2 a 4 controlan además por sector de
actividad. Errores estándar entre paréntesis. *** p \textless{} 0,001;
** p \textless{} 0,01; * p \textless{} 0,05.

{

\begin{longtable}[]{@{}
  >{\raggedright\arraybackslash}p{(\linewidth - 8\tabcolsep) * \real{0.3684}}
  >{\centering\arraybackslash}p{(\linewidth - 8\tabcolsep) * \real{0.1579}}
  >{\centering\arraybackslash}p{(\linewidth - 8\tabcolsep) * \real{0.1579}}
  >{\centering\arraybackslash}p{(\linewidth - 8\tabcolsep) * \real{0.1579}}
  >{\centering\arraybackslash}p{(\linewidth - 8\tabcolsep) * \real{0.1579}}@{}}

\caption{\label{tbl-a-proteccion}Modelos secuenciales de la dirección:
protección}

\tabularnewline

\toprule\noalign{}
\begin{minipage}[b]{\linewidth}\raggedright
\end{minipage} & \begin{minipage}[b]{\linewidth}\centering
Modelo 1
\end{minipage} & \begin{minipage}[b]{\linewidth}\centering
Modelo 2
\end{minipage} & \begin{minipage}[b]{\linewidth}\centering
Modelo 3
\end{minipage} & \begin{minipage}[b]{\linewidth}\centering
Modelo 4
\end{minipage} \\
\midrule\noalign{}
\endhead
\bottomrule\noalign{}
\endlastfoot
Intercepto & 0,061 *** & 0,100 *** & 0,081 ** & 0,082 ** \\
& (0,012) & (0,030) & (0,031) & (0,031) \\
\textbf{Nivel individual} & & & & \\
Amenaza tecnológica & 0,016 *** & 0,011 * & 0,009 & 0,009 \\
& (0,005) & (0,005) & (0,005) & (0,005) \\
Beneficio esperado & -0,051 *** & -0,044 *** & -0,044 *** & -0,044
*** \\
& (0,005) & (0,005) & (0,005) & (0,005) \\
Expectativas negativas sobre la IA & 0,112 *** & 0,109 *** & 0,108 *** &
0,108 *** \\
& (0,005) & (0,005) & (0,005) & (0,005) \\
Expectativas positivas sobre la IA & -0,088 *** & -0,083 *** & -0,082
*** & -0,082 *** \\
& (0,006) & (0,006) & (0,006) & (0,006) \\
Preocupación por perder el empleo & & & 0,008 * & 0,007 * \\
& & & (0,003) & (0,003) \\
Mujer & & 0,022 ** & 0,020 ** & 0,020 ** \\
& & (0,007) & (0,007) & (0,007) \\
25 a 34 años & & 0,006 & 0,005 & 0,005 \\
& & (0,012) & (0,012) & (0,012) \\
35 a 44 años & & 0,009 & 0,008 & 0,008 \\
& & (0,012) & (0,012) & (0,012) \\
45 a 54 años & & 0,029 * & 0,028 * & 0,028 * \\
& & (0,012) & (0,012) & (0,012) \\
55 a 64 años & & 0,015 & 0,015 & 0,015 \\
& & (0,013) & (0,013) & (0,013) \\
Educación media & & -0,009 & -0,009 & -0,009 \\
& & (0,013) & (0,013) & (0,013) \\
Educación alta & & -0,101 *** & -0,101 *** & -0,101 *** \\
& & (0,014) & (0,014) & (0,014) \\
Decil de ingreso & & -0,002 & -0,001 & -0,001 \\
& & (0,001) & (0,001) & (0,001) \\
Asalariado temporal o sin contrato & & -0,018 & -0,019 * & -0,019 * \\
& & (0,010) & (0,010) & (0,010) \\
Independiente & & -0,027 ** & -0,028 ** & -0,028 ** \\
& & (0,010) & (0,010) & (0,010) \\
Trabajo en plataformas & & 0,013 & 0,012 & 0,012 \\
& & (0,011) & (0,011) & (0,011) \\
Sector público & & 0,008 & 0,009 & 0,009 \\
& & (0,008) & (0,008) & (0,008) \\
\textbf{Nivel país} & & & & \\
PIB per cápita (log) & & & & -0,010 \\
& & & & (0,017) \\
Gini & & & & 0,004 \\
& & & & (0,017) \\
Prestaciones por desempleo & & & & 0,001 \\
& & & & (0,016) \\
Políticas de mercado laboral & & & & -0,014 \\
& & & & (0,017) \\
Log verosimilitud & -13173,919 & -13090,977 & -13093,215 & -13104,640 \\
N & 19.731 & 19.731 & 19.731 & 19.731 \\
Países & 27 & 27 & 27 & 27 \\
Varianza del intercepto (país) & 0,004 & 0,004 & 0,004 & 0,004 \\
Varianza residual & 0,199 & 0,196 & 0,196 & 0,196 \\

\end{longtable}

}

\emph{Nota}: Categorías de referencia: hombre, 18 a 24 años, educación
baja, asalariado con contrato indefinido, sin trabajo en plataformas y
no empleado en el sector público (asalariados del sector privado e
independientes). Los modelos 2 a 4 controlan además por sector de
actividad. Errores estándar entre paréntesis. *** p \textless{} 0,001;
** p \textless{} 0,01; * p \textless{} 0,05.

{

\begin{longtable}[]{@{}
  >{\raggedright\arraybackslash}p{(\linewidth - 4\tabcolsep) * \real{0.4605}}
  >{\centering\arraybackslash}p{(\linewidth - 4\tabcolsep) * \real{0.1974}}
  >{\centering\arraybackslash}p{(\linewidth - 4\tabcolsep) * \real{0.3421}}@{}}

\caption{\label{tbl-a-sin-subsidios}Dirección de la protección con y sin
subsidios}

\tabularnewline

\toprule\noalign{}
\begin{minipage}[b]{\linewidth}\raggedright
\end{minipage} & \begin{minipage}[b]{\linewidth}\centering
Protección
\end{minipage} & \begin{minipage}[b]{\linewidth}\centering
Protección sin subsidios
\end{minipage} \\
\midrule\noalign{}
\endhead
\bottomrule\noalign{}
\endlastfoot
Intercepto & 0,082 ** & 0,187 *** \\
& (0,031) & (0,039) \\
\textbf{Nivel individual} & & \\
Amenaza tecnológica & 0,009 & 0,022 *** \\
& (0,005) & (0,006) \\
Beneficio esperado & -0,044 *** & -0,057 *** \\
& (0,005) & (0,006) \\
Expectativas negativas sobre la IA & 0,108 *** & 0,132 *** \\
& (0,005) & (0,006) \\
Expectativas positivas sobre la IA & -0,082 *** & -0,108 *** \\
& (0,006) & (0,007) \\
Preocupación por perder el empleo & 0,007 * & 0,002 \\
& (0,003) & (0,004) \\
Mujer & 0,020 ** & -0,010 \\
& (0,007) & (0,008) \\
25 a 34 años & 0,005 & -0,011 \\
& (0,012) & (0,015) \\
35 a 44 años & 0,008 & 0,001 \\
& (0,012) & (0,015) \\
45 a 54 años & 0,028 * & 0,023 \\
& (0,012) & (0,015) \\
55 a 64 años & 0,015 & 0,006 \\
& (0,013) & (0,016) \\
Educación media & -0,009 & -0,016 \\
& (0,013) & (0,016) \\
Educación alta & -0,101 *** & -0,106 *** \\
& (0,014) & (0,017) \\
Decil de ingreso & -0,001 & -0,001 \\
& (0,001) & (0,001) \\
Asalariado temporal o sin contrato & -0,019 * & -0,019 \\
& (0,010) & (0,012) \\
Independiente & -0,028 ** & -0,030 * \\
& (0,010) & (0,013) \\
Trabajo en plataformas & 0,012 & 0,010 \\
& (0,011) & (0,014) \\
Sector público & 0,009 & 0,022 * \\
& (0,008) & (0,010) \\
\textbf{Nivel país} & & \\
PIB per cápita (log) & -0,010 & -0,011 \\
& (0,017) & (0,024) \\
Gini & 0,004 & 0,006 \\
& (0,017) & (0,023) \\
Prestaciones por desempleo & 0,001 & -0,004 \\
& (0,016) & (0,023) \\
Políticas de mercado laboral & -0,014 & -0,012 \\
& (0,017) & (0,024) \\
Log verosimilitud & -13104,640 & -17068,172 \\
N & 19.731 & 19.708 \\
Países & 27 & 27 \\
Varianza del intercepto (país) & 0,004 & 0,008 \\
Varianza residual & 0,196 & 0,293 \\

\end{longtable}

}

\emph{Nota}: Categorías de referencia: hombre, 18 a 24 años, educación
baja, asalariado con contrato indefinido, sin trabajo en plataformas y
no empleado en el sector público (asalariados del sector privado e
independientes). Todos los modelos controlan además por sector de
actividad. Errores estándar entre paréntesis. *** p \textless{} 0,001;
** p \textless{} 0,01; * p \textless{} 0,05. Los indicadores de nivel
país están estandarizados entre los 27 países. PIB per cápita en paridad
de poder adquisitivo, en logaritmo, y coeficiente de Gini, ambos del
Banco Mundial. Prestaciones por desempleo: porcentaje de personas
desempleadas que reciben prestaciones monetarias por desempleo (OIT,
indicador 1.3.1 de los ODS). Políticas de mercado laboral: gasto público
en políticas activas del mercado laboral como porcentaje del PIB,
dividido por la tasa de desempleo (OCDE, SOCX).

{

\begin{longtable}[]{@{}
  >{\raggedright\arraybackslash}p{(\linewidth - 8\tabcolsep) * \real{0.3241}}
  >{\centering\arraybackslash}p{(\linewidth - 8\tabcolsep) * \real{0.1759}}
  >{\centering\arraybackslash}p{(\linewidth - 8\tabcolsep) * \real{0.1389}}
  >{\centering\arraybackslash}p{(\linewidth - 8\tabcolsep) * \real{0.2222}}
  >{\centering\arraybackslash}p{(\linewidth - 8\tabcolsep) * \real{0.1389}}@{}}

\caption{\label{tbl-a-items}Dirección de cada medida de protección}

\tabularnewline

\toprule\noalign{}
\begin{minipage}[b]{\linewidth}\raggedright
\end{minipage} & \begin{minipage}[b]{\linewidth}\centering
Impuesto a robots
\end{minipage} & \begin{minipage}[b]{\linewidth}\centering
Subsidios
\end{minipage} & \begin{minipage}[b]{\linewidth}\centering
Prioridad a nacionales
\end{minipage} & \begin{minipage}[b]{\linewidth}\centering
Aranceles
\end{minipage} \\
\midrule\noalign{}
\endhead
\bottomrule\noalign{}
\endlastfoot
Intercepto & 0,126 * & -0,229 *** & 0,367 *** & 0,065 \\
& (0,063) & (0,052) & (0,067) & (0,064) \\
\textbf{Nivel individual} & & & & \\
Amenaza tecnológica & 0,118 *** & -0,033 *** & -0,103 *** & 0,029 ** \\
& (0,010) & (0,009) & (0,010) & (0,009) \\
Beneficio esperado & -0,080 *** & -0,006 & -0,042 *** & -0,023 * \\
& (0,011) & (0,010) & (0,011) & (0,010) \\
Expectativas negativas sobre la IA & 0,205 *** & 0,028 ** & 0,110 *** &
0,068 *** \\
& (0,011) & (0,010) & (0,011) & (0,011) \\
Expectativas positivas sobre la IA & -0,154 *** & 0,001 & -0,136 *** &
-0,033 ** \\
& (0,012) & (0,010) & (0,012) & (0,011) \\
Preocupación por perder el empleo & 0,003 & 0,026 *** & 0,004 &
-0,004 \\
& (0,007) & (0,006) & (0,007) & (0,007) \\
Mujer & -0,079 *** & 0,109 *** & -0,022 & 0,051 *** \\
& (0,014) & (0,012) & (0,014) & (0,013) \\
25 a 34 años & -0,090 *** & 0,057 * & 0,032 & 0,022 \\
& (0,026) & (0,022) & (0,026) & (0,025) \\
35 a 44 años & -0,098 *** & 0,039 & 0,049 & 0,037 \\
& (0,026) & (0,022) & (0,026) & (0,024) \\
45 a 54 años & -0,099 *** & 0,050 * & 0,099 *** & 0,042 \\
& (0,026) & (0,023) & (0,027) & (0,025) \\
55 a 64 años & -0,102 *** & 0,046 & 0,088 ** & -0,010 \\
& (0,027) & (0,024) & (0,028) & (0,026) \\
Educación media & 0,007 & 0,014 & 0,000 & -0,062 * \\
& (0,028) & (0,024) & (0,028) & (0,026) \\
Educación alta & -0,040 & -0,075 ** & -0,137 *** & -0,138 *** \\
& (0,028) & (0,025) & (0,029) & (0,027) \\
Decil de ingreso & -0,000 & -0,004 & 0,001 & -0,004 \\
& (0,002) & (0,002) & (0,002) & (0,002) \\
Asalariado temporal o sin contrato & -0,029 & -0,025 & -0,001 &
-0,004 \\
& (0,020) & (0,018) & (0,021) & (0,019) \\
Independiente & -0,057 ** & -0,028 & -0,008 & -0,025 \\
& (0,021) & (0,018) & (0,022) & (0,020) \\
Trabajo en plataformas & 0,024 & 0,002 & -0,040 & 0,052 * \\
& (0,024) & (0,021) & (0,024) & (0,023) \\
Sector público & 0,004 & -0,027 & 0,012 & 0,027 \\
& (0,017) & (0,015) & (0,017) & (0,016) \\
\textbf{Nivel país} & & & & \\
PIB per cápita (log) & 0,015 & -0,007 & -0,027 & -0,025 \\
& (0,031) & (0,015) & (0,039) & (0,042) \\
Gini & 0,006 & -0,006 & 0,011 & -0,001 \\
& (0,030) & (0,015) & (0,038) & (0,041) \\
Prestaciones por desempleo & 0,043 & 0,014 & -0,021 & -0,020 \\
& (0,029) & (0,014) & (0,037) & (0,040) \\
Políticas de mercado laboral & -0,026 & -0,023 & 0,004 & -0,012 \\
& (0,031) & (0,015) & (0,039) & (0,042) \\
Log verosimilitud & -24327,606 & -21888,612 & -24841,400 & -23500,264 \\
N & 17.987 & 17.987 & 17.987 & 17.987 \\
Países & 27 & 27 & 27 & 27 \\
Varianza del intercepto (país) & 0,012 & 0,002 & 0,020 & 0,024 \\
Varianza residual & 0,775 & 0,592 & 0,821 & 0,706 \\

\end{longtable}

}

\emph{Nota}: Categorías de referencia: hombre, 18 a 24 años, educación
baja, asalariado con contrato indefinido, sin trabajo en plataformas y
no empleado en el sector público (asalariados del sector privado e
independientes). Todos los modelos controlan además por sector de
actividad. Errores estándar entre paréntesis. *** p \textless{} 0,001;
** p \textless{} 0,01; * p \textless{} 0,05. Los indicadores de nivel
país están estandarizados entre los 27 países. PIB per cápita en paridad
de poder adquisitivo, en logaritmo, y coeficiente de Gini, ambos del
Banco Mundial. Prestaciones por desempleo: porcentaje de personas
desempleadas que reciben prestaciones monetarias por desempleo (OIT,
indicador 1.3.1 de los ODS). Políticas de mercado laboral: gasto público
en políticas activas del mercado laboral como porcentaje del PIB,
dividido por la tasa de desempleo (OCDE, SOCX).

{

\begin{longtable}[]{@{}
  >{\raggedright\arraybackslash}p{(\linewidth - 8\tabcolsep) * \real{0.4286}}
  >{\centering\arraybackslash}p{(\linewidth - 8\tabcolsep) * \real{0.1429}}
  >{\centering\arraybackslash}p{(\linewidth - 8\tabcolsep) * \real{0.1429}}
  >{\centering\arraybackslash}p{(\linewidth - 8\tabcolsep) * \real{0.1429}}
  >{\centering\arraybackslash}p{(\linewidth - 8\tabcolsep) * \real{0.1429}}@{}}

\caption{\label{tbl-a-moderacion-edad}Interacciones entre la amenaza
tecnológica y el grupo de edad}

\tabularnewline

\toprule\noalign{}
\begin{minipage}[b]{\linewidth}\raggedright
\end{minipage} & \begin{minipage}[b]{\linewidth}\centering
Volumen
\end{minipage} & \begin{minipage}[b]{\linewidth}\centering
Inversión
\end{minipage} & \begin{minipage}[b]{\linewidth}\centering
Compensación
\end{minipage} & \begin{minipage}[b]{\linewidth}\centering
Protección
\end{minipage} \\
\midrule\noalign{}
\endhead
\bottomrule\noalign{}
\endlastfoot
Intercepto & 3,297 *** & 0,149 *** & -0,129 ** & 0,081 ** \\
& (0,035) & (0,035) & (0,046) & (0,031) \\
\textbf{Nivel individual} & & & & \\
Amenaza tecnológica & 0,019 & -0,119 *** & -0,110 *** & 0,014 \\
& (0,014) & (0,013) & (0,019) & (0,012) \\
Beneficio esperado & 0,091 *** & 0,058 *** & 0,023 ** & -0,044 *** \\
& (0,006) & (0,005) & (0,008) & (0,005) \\
Expectativas negativas sobre la IA & 0,192 *** & -0,011 * & 0,059 *** &
0,108 *** \\
& (0,006) & (0,006) & (0,008) & (0,005) \\
Expectativas positivas sobre la IA & 0,142 *** & -0,004 & -0,044 *** &
-0,082 *** \\
& (0,006) & (0,006) & (0,008) & (0,006) \\
Preocupación por perder el empleo & 0,042 *** & -0,017 *** & 0,050 *** &
0,007 * \\
& (0,004) & (0,004) & (0,005) & (0,003) \\
Mujer & 0,025 *** & -0,016 * & 0,035 *** & 0,021 ** \\
& (0,007) & (0,007) & (0,010) & (0,007) \\
25 a 34 años & 0,068 *** & 0,047 *** & 0,020 & 0,006 \\
& (0,014) & (0,013) & (0,019) & (0,013) \\
35 a 44 años & 0,068 *** & 0,096 *** & 0,019 & 0,009 \\
& (0,014) & (0,013) & (0,019) & (0,013) \\
45 a 54 años & 0,068 *** & 0,148 *** & 0,020 & 0,029 * \\
& (0,014) & (0,013) & (0,019) & (0,013) \\
55 a 64 años & 0,075 *** & 0,183 *** & 0,037 & 0,015 \\
& (0,015) & (0,014) & (0,020) & (0,013) \\
Educación media & -0,009 & 0,070 *** & -0,015 & -0,010 \\
& (0,015) & (0,014) & (0,020) & (0,013) \\
Educación alta & 0,009 & 0,133 *** & -0,072 *** & -0,101 *** \\
& (0,015) & (0,014) & (0,021) & (0,014) \\
Decil de ingreso & -0,002 & 0,008 *** & -0,016 *** & -0,001 \\
& (0,001) & (0,001) & (0,002) & (0,001) \\
Asalariado temporal o sin contrato & 0,012 & 0,027 ** & 0,001 & -0,019
* \\
& (0,011) & (0,010) & (0,015) & (0,010) \\
Independiente & -0,066 *** & 0,012 & 0,003 & -0,029 ** \\
& (0,011) & (0,011) & (0,015) & (0,010) \\
Trabajo en plataformas & 0,122 *** & -0,102 *** & -0,040 * & 0,012 \\
& (0,013) & (0,012) & (0,017) & (0,011) \\
Sector público & -0,026 ** & -0,049 *** & 0,007 & 0,009 \\
& (0,009) & (0,009) & (0,012) & (0,008) \\
\textbf{Nivel país} & & & & \\
PIB per cápita (log) & 0,005 & -0,003 & 0,008 & -0,010 \\
& (0,022) & (0,025) & (0,023) & (0,017) \\
Gini & 0,037 & 0,029 & -0,012 & 0,004 \\
& (0,022) & (0,024) & (0,022) & (0,017) \\
Prestaciones por desempleo & -0,024 & 0,010 & -0,043 * & 0,001 \\
& (0,021) & (0,023) & (0,022) & (0,016) \\
Políticas de mercado laboral & -0,035 & 0,025 & -0,038 & -0,013 \\
& (0,022) & (0,025) & (0,023) & (0,017) \\
\textbf{Interacciones con variables individuales} & & & & \\
Amenaza tecnológica × 25 a 34 años & -0,022 & -0,029 & 0,034 & -0,001 \\
& (0,016) & (0,015) & (0,022) & (0,014) \\
Amenaza tecnológica × 35 a 44 años & -0,005 & -0,049 ** & 0,076 *** &
-0,011 \\
& (0,016) & (0,015) & (0,022) & (0,014) \\
Amenaza tecnológica × 45 a 54 años & -0,022 & -0,055 *** & 0,084 *** &
0,001 \\
& (0,016) & (0,015) & (0,022) & (0,015) \\
Amenaza tecnológica × 55 a 64 años & -0,009 & -0,034 * & 0,132 *** &
-0,009 \\
& (0,017) & (0,017) & (0,024) & (0,016) \\
Log verosimilitud & -15092,223 & -14153,749 & -21259,222 & -13118,115 \\
N & 19.731 & 19.731 & 19.731 & 19.731 \\
Países & 27 & 27 & 27 & 27 \\
Varianza del intercepto (país) & 0,007 & 0,008 & 0,007 & 0,004 \\
Varianza residual & 0,239 & 0,217 & 0,447 & 0,196 \\

\end{longtable}

}

\emph{Nota}: Categorías de referencia: hombre, 18 a 24 años, educación
baja, asalariado con contrato indefinido, sin trabajo en plataformas y
no empleado en el sector público (asalariados del sector privado e
independientes). Todos los modelos controlan además por sector de
actividad. Errores estándar entre paréntesis. *** p \textless{} 0,001;
** p \textless{} 0,01; * p \textless{} 0,05. Los indicadores de nivel
país están estandarizados entre los 27 países. PIB per cápita en paridad
de poder adquisitivo, en logaritmo, y coeficiente de Gini, ambos del
Banco Mundial. Prestaciones por desempleo: porcentaje de personas
desempleadas que reciben prestaciones monetarias por desempleo (OIT,
indicador 1.3.1 de los ODS). Políticas de mercado laboral: gasto público
en políticas activas del mercado laboral como porcentaje del PIB,
dividido por la tasa de desempleo (OCDE, SOCX).

{

\begin{longtable}[]{@{}
  >{\raggedright\arraybackslash}p{(\linewidth - 8\tabcolsep) * \real{0.4286}}
  >{\centering\arraybackslash}p{(\linewidth - 8\tabcolsep) * \real{0.1429}}
  >{\centering\arraybackslash}p{(\linewidth - 8\tabcolsep) * \real{0.1429}}
  >{\centering\arraybackslash}p{(\linewidth - 8\tabcolsep) * \real{0.1429}}
  >{\centering\arraybackslash}p{(\linewidth - 8\tabcolsep) * \real{0.1429}}@{}}

\caption{\label{tbl-a-moderacion-educacion}Interacciones entre la
amenaza tecnológica y el nivel educacional}

\tabularnewline

\toprule\noalign{}
\begin{minipage}[b]{\linewidth}\raggedright
\end{minipage} & \begin{minipage}[b]{\linewidth}\centering
Volumen
\end{minipage} & \begin{minipage}[b]{\linewidth}\centering
Inversión
\end{minipage} & \begin{minipage}[b]{\linewidth}\centering
Compensación
\end{minipage} & \begin{minipage}[b]{\linewidth}\centering
Protección
\end{minipage} \\
\midrule\noalign{}
\endhead
\bottomrule\noalign{}
\endlastfoot
Intercepto & 3,307 *** & 0,152 *** & -0,146 ** & 0,090 ** \\
& (0,035) & (0,035) & (0,046) & (0,031) \\
\textbf{Nivel individual} & & & & \\
Amenaza tecnológica & -0,015 & -0,086 *** & -0,093 *** & -0,039 ** \\
& (0,016) & (0,015) & (0,022) & (0,014) \\
Beneficio esperado & 0,093 *** & 0,056 *** & 0,025 ** & -0,042 *** \\
& (0,006) & (0,005) & (0,008) & (0,005) \\
Expectativas negativas sobre la IA & 0,190 *** & -0,009 & 0,057 *** &
0,106 *** \\
& (0,006) & (0,006) & (0,008) & (0,005) \\
Expectativas positivas sobre la IA & 0,141 *** & -0,003 & -0,045 *** &
-0,083 *** \\
& (0,006) & (0,006) & (0,008) & (0,006) \\
Preocupación por perder el empleo & 0,041 *** & -0,016 *** & 0,050 *** &
0,006 \\
& (0,004) & (0,004) & (0,005) & (0,003) \\
Mujer & 0,027 *** & -0,020 ** & 0,040 *** & 0,023 *** \\
& (0,007) & (0,007) & (0,010) & (0,007) \\
25 a 34 años & 0,063 *** & 0,038 ** & 0,035 & 0,005 \\
& (0,014) & (0,013) & (0,019) & (0,012) \\
35 a 44 años & 0,063 *** & 0,087 *** & 0,036 & 0,006 \\
& (0,013) & (0,013) & (0,018) & (0,012) \\
45 a 54 años & 0,065 *** & 0,140 *** & 0,036 & 0,028 * \\
& (0,014) & (0,013) & (0,019) & (0,012) \\
55 a 64 años & 0,070 *** & 0,173 *** & 0,040 * & 0,014 \\
& (0,014) & (0,014) & (0,020) & (0,013) \\
Educación media & -0,011 & 0,072 *** & -0,016 & -0,012 \\
& (0,015) & (0,014) & (0,020) & (0,013) \\
Educación alta & 0,008 & 0,136 *** & -0,074 *** & -0,103 *** \\
& (0,015) & (0,014) & (0,021) & (0,014) \\
Decil de ingreso & -0,002 & 0,008 *** & -0,016 *** & -0,002 \\
& (0,001) & (0,001) & (0,002) & (0,001) \\
Asalariado temporal o sin contrato & 0,013 & 0,026 ** & 0,002 &
-0,018 \\
& (0,011) & (0,010) & (0,015) & (0,010) \\
Independiente & -0,065 *** & 0,010 & 0,002 & -0,026 ** \\
& (0,011) & (0,011) & (0,015) & (0,010) \\
Trabajo en plataformas & 0,119 *** & -0,099 *** & -0,046 ** & 0,009 \\
& (0,013) & (0,012) & (0,017) & (0,011) \\
Sector público & -0,024 ** & -0,049 *** & 0,005 & 0,011 \\
& (0,009) & (0,009) & (0,012) & (0,008) \\
\textbf{Nivel país} & & & & \\
PIB per cápita (log) & 0,005 & -0,003 & 0,008 & -0,010 \\
& (0,022) & (0,025) & (0,023) & (0,017) \\
Gini & 0,038 & 0,028 & -0,011 & 0,005 \\
& (0,021) & (0,024) & (0,022) & (0,017) \\
Prestaciones por desempleo & -0,024 & 0,011 & -0,044 * & -0,000 \\
& (0,021) & (0,024) & (0,022) & (0,016) \\
Políticas de mercado laboral & -0,035 & 0,025 & -0,038 & -0,014 \\
& (0,022) & (0,025) & (0,023) & (0,017) \\
\textbf{Interacciones con variables individuales} & & & & \\
Amenaza tecnológica × Educación media & -0,013 & -0,050 ** & 0,043 &
0,020 \\
& (0,017) & (0,016) & (0,023) & (0,015) \\
Amenaza tecnológica × Educación alta & 0,055 *** & -0,102 *** & 0,065 **
& 0,082 *** \\
& (0,017) & (0,016) & (0,023) & (0,015) \\
Log verosimilitud & -15056,864 & -14123,394 & -21268,913 & -13075,727 \\
N & 19.731 & 19.731 & 19.731 & 19.731 \\
Países & 27 & 27 & 27 & 27 \\
Varianza del intercepto (país) & 0,007 & 0,009 & 0,007 & 0,004 \\
Varianza residual & 0,238 & 0,217 & 0,448 & 0,195 \\

\end{longtable}

}

\emph{Nota}: Categorías de referencia: hombre, 18 a 24 años, educación
baja, asalariado con contrato indefinido, sin trabajo en plataformas y
no empleado en el sector público (asalariados del sector privado e
independientes). Todos los modelos controlan además por sector de
actividad. Errores estándar entre paréntesis. *** p \textless{} 0,001;
** p \textless{} 0,01; * p \textless{} 0,05. Los indicadores de nivel
país están estandarizados entre los 27 países. PIB per cápita en paridad
de poder adquisitivo, en logaritmo, y coeficiente de Gini, ambos del
Banco Mundial. Prestaciones por desempleo: porcentaje de personas
desempleadas que reciben prestaciones monetarias por desempleo (OIT,
indicador 1.3.1 de los ODS). Políticas de mercado laboral: gasto público
en políticas activas del mercado laboral como porcentaje del PIB,
dividido por la tasa de desempleo (OCDE, SOCX).

{

\begin{longtable}[]{@{}
  >{\raggedright\arraybackslash}p{(\linewidth - 8\tabcolsep) * \real{0.1549}}
  >{\centering\arraybackslash}p{(\linewidth - 8\tabcolsep) * \real{0.2113}}
  >{\centering\arraybackslash}p{(\linewidth - 8\tabcolsep) * \real{0.2113}}
  >{\centering\arraybackslash}p{(\linewidth - 8\tabcolsep) * \real{0.2113}}
  >{\centering\arraybackslash}p{(\linewidth - 8\tabcolsep) * \real{0.2113}}@{}}

\caption{\label{tbl-a-moderacion-pendientes}Prueba conjunta y pendiente
de la amenaza tecnológica por grupo de edad y nivel educacional}

\tabularnewline

\toprule\noalign{}
\begin{minipage}[b]{\linewidth}\raggedright
\end{minipage} & \begin{minipage}[b]{\linewidth}\centering
Volumen
\end{minipage} & \begin{minipage}[b]{\linewidth}\centering
Inversión
\end{minipage} & \begin{minipage}[b]{\linewidth}\centering
Compensación
\end{minipage} & \begin{minipage}[b]{\linewidth}\centering
Protección
\end{minipage} \\
\midrule\noalign{}
\endhead
\bottomrule\noalign{}
\endlastfoot
\textbf{Grupo de edad} & & & & \\
18 a 24 años & 0,019 (0,014) & -0,119 (0,013)*** & -0,110 (0,019)*** &
0,014 (0,012) \\
25 a 34 años & -0,003 (0,009) & -0,148 (0,009)*** & -0,076 (0,012)*** &
0,012 (0,008) \\
35 a 44 años & 0,014 (0,009) & -0,168 (0,008)*** & -0,034 (0,012)** &
0,002 (0,008) \\
45 a 54 años & -0,004 (0,010) & -0,174 (0,009)*** & -0,026 (0,013)* &
0,014 (0,009) \\
55 a 64 años & 0,009 (0,012) & -0,153 (0,011)*** & 0,021 (0,016) & 0,004
(0,010) \\
Prueba conjunta & χ² = 4,12 (4 gl), p = 0,390 & χ² = 15,96 (4 gl), p =
0,003 & χ² = 41,27 (4 gl), p = 0,000 & χ² = 1,71 (4 gl), p = 0,789 \\
\textbf{Nivel educacional} & & & & \\
Educación baja & -0,015 (0,016) & -0,086 (0,015)*** & -0,093 (0,022)***
& -0,039 (0,014)** \\
Educación media & -0,028 (0,007)*** & -0,137 (0,007)*** & -0,049
(0,010)*** & -0,019 (0,006)** \\
Educación alta & 0,041 (0,007)*** & -0,188 (0,007)*** & -0,028 (0,009)**
& 0,043 (0,006)*** \\
Prueba conjunta & χ² = 61,11 (2 gl), p = 0,000 & χ² = 62,69 (2 gl), p =
0,000 & χ² = 9,25 (2 gl), p = 0,010 & χ² = 72,38 (2 gl), p = 0,000 \\

\end{longtable}

}

\emph{Nota}: pendiente de la amenaza tecnológica en cada categoría,
calculada como la suma del coeficiente de la amenaza y de su interacción
con la categoría, con el error estándar de esa suma. Prueba conjunta:
razón de verosimilitud de cada modelo frente al modelo 4 sin la
interacción. Categorías de referencia: hombre, 18 a 24 años, educación
baja, asalariado con contrato indefinido, sin trabajo en plataformas y
no empleado en el sector público (asalariados del sector privado e
independientes). Todos los modelos controlan además por sector de
actividad. Errores estándar entre paréntesis. *** p \textless{} 0,001;
** p \textless{} 0,01; * p \textless{} 0,05. Los indicadores de nivel
país están estandarizados entre los 27 países. PIB per cápita en paridad
de poder adquisitivo, en logaritmo, y coeficiente de Gini, ambos del
Banco Mundial. Prestaciones por desempleo: porcentaje de personas
desempleadas que reciben prestaciones monetarias por desempleo (OIT,
indicador 1.3.1 de los ODS). Políticas de mercado laboral: gasto público
en políticas activas del mercado laboral como porcentaje del PIB,
dividido por la tasa de desempleo (OCDE, SOCX).

{

\begin{longtable}[]{@{}
  >{\raggedright\arraybackslash}p{(\linewidth - 10\tabcolsep) * \real{0.2500}}
  >{\raggedleft\arraybackslash}p{(\linewidth - 10\tabcolsep) * \real{0.0357}}
  >{\centering\arraybackslash}p{(\linewidth - 10\tabcolsep) * \real{0.1786}}
  >{\centering\arraybackslash}p{(\linewidth - 10\tabcolsep) * \real{0.1786}}
  >{\centering\arraybackslash}p{(\linewidth - 10\tabcolsep) * \real{0.1786}}
  >{\centering\arraybackslash}p{(\linewidth - 10\tabcolsep) * \real{0.1786}}@{}}

\caption{\label{tbl-a-sector}Pendiente de la amenaza tecnológica por
sector de actividad (análisis exploratorio)}

\tabularnewline

\toprule\noalign{}
\begin{minipage}[b]{\linewidth}\raggedright
Sector
\end{minipage} & \begin{minipage}[b]{\linewidth}\raggedleft
N
\end{minipage} & \begin{minipage}[b]{\linewidth}\centering
Volumen
\end{minipage} & \begin{minipage}[b]{\linewidth}\centering
Inversión
\end{minipage} & \begin{minipage}[b]{\linewidth}\centering
Compensación
\end{minipage} & \begin{minipage}[b]{\linewidth}\centering
Protección
\end{minipage} \\
\midrule\noalign{}
\endhead
\bottomrule\noalign{}
\endlastfoot
Agricultura y minería & 445 & 0,060 (0,027)* & -0,121 (0,026)*** &
-0,033 (0,037) & -0,039 (0,024) \\
Manufactura y servicios básicos & 2.264 & 0,012 (0,013) & -0,165
(0,013)*** & -0,084 (0,018)*** & 0,002 (0,012) \\
Construcción & 1.167 & 0,006 (0,017) & -0,126 (0,016)*** & -0,061
(0,023)** & -0,032 (0,015)* \\
Comercio, transporte y hotelería & 2.831 & -0,013 (0,012) & -0,144
(0,011)*** & -0,050 (0,016)** & -0,023 (0,010)* \\
TIC, finanzas y servicios a empresas & 4.598 & 0,066 (0,009)*** & -0,181
(0,009)*** & -0,022 (0,013) & 0,040 (0,008)*** \\
Administración pública, educación y salud & 4.474 & -0,043 (0,010)*** &
-0,171 (0,010)*** & -0,020 (0,014) & 0,020 (0,009)* \\
Otros servicios & 2.878 & -0,004 (0,011) & -0,130 (0,011)*** & -0,061
(0,016)*** & 0,007 (0,010) \\
Otro o no declarado & 1.074 & -0,041 (0,018)* & -0,150 (0,017)*** &
-0,040 (0,025) & 0,007 (0,016) \\
Prueba conjunta & & χ² = 86,70 (7 gl), p = 0,000 & χ² = 25,36 (7 gl), p
= 0,001 & χ² = 14,10 (7 gl), p = 0,050 & χ² = 37,74 (7 gl), p = 0,000 \\

\end{longtable}

}

\emph{Nota}: modelos con una interacción entre la amenaza tecnológica y
el sector de actividad. La pendiente de cada sector es la suma del
coeficiente de la amenaza y de su interacción con el sector, con el
error estándar de esa suma. Prueba conjunta: razón de verosimilitud
frente al modelo 4 sin la interacción. N: personas de cada sector en la
muestra analítica. Categorías de referencia: hombre, 18 a 24 años,
educación baja, asalariado con contrato indefinido, sin trabajo en
plataformas y no empleado en el sector público (asalariados del sector
privado e independientes). Todos los modelos controlan además por sector
de actividad. Errores estándar entre paréntesis. *** p \textless{}
0,001; ** p \textless{} 0,01; * p \textless{} 0,05. Los indicadores de
nivel país están estandarizados entre los 27 países. PIB per cápita en
paridad de poder adquisitivo, en logaritmo, y coeficiente de Gini, ambos
del Banco Mundial. Prestaciones por desempleo: porcentaje de personas
desempleadas que reciben prestaciones monetarias por desempleo (OIT,
indicador 1.3.1 de los ODS). Políticas de mercado laboral: gasto público
en políticas activas del mercado laboral como porcentaje del PIB,
dividido por la tasa de desempleo (OCDE, SOCX).

{

\begin{longtable}[]{@{}
  >{\raggedright\arraybackslash}p{(\linewidth - 6\tabcolsep) * \real{0.3248}}
  >{\centering\arraybackslash}p{(\linewidth - 6\tabcolsep) * \real{0.1592}}
  >{\centering\arraybackslash}p{(\linewidth - 6\tabcolsep) * \real{0.2484}}
  >{\centering\arraybackslash}p{(\linewidth - 6\tabcolsep) * \real{0.2675}}@{}}

\caption{\label{tbl-a-gasto-social}Prestaciones por desempleo y gasto
social público}

\tabularnewline

\toprule\noalign{}
\begin{minipage}[b]{\linewidth}\raggedright
\end{minipage} & \begin{minipage}[b]{\linewidth}\centering
Modelo 4 + gasto social
\end{minipage} & \begin{minipage}[b]{\linewidth}\centering
Gasto social en lugar de prestaciones
\end{minipage} & \begin{minipage}[b]{\linewidth}\centering
Amenaza × prestaciones, con gasto social
\end{minipage} \\
\midrule\noalign{}
\endhead
\bottomrule\noalign{}
\endlastfoot
Intercepto & -0,151 *** & -0,151 *** & -0,147 ** \\
& (0,045) & (0,045) & (0,045) \\
\textbf{Nivel individual} & & & \\
Amenaza tecnológica & -0,042 *** & -0,042 *** & -0,041 *** \\
& (0,007) & (0,007) & (0,012) \\
Beneficio esperado & 0,024 ** & 0,024 ** & 0,022 ** \\
& (0,008) & (0,008) & (0,008) \\
Expectativas negativas sobre la IA & 0,058 *** & 0,058 *** & 0,058
*** \\
& (0,008) & (0,008) & (0,008) \\
Expectativas positivas sobre la IA & -0,045 *** & -0,045 *** & -0,044
*** \\
& (0,008) & (0,008) & (0,008) \\
Preocupación por perder el empleo & 0,050 *** & 0,050 *** & 0,049 *** \\
& (0,005) & (0,005) & (0,005) \\
Mujer & 0,039 *** & 0,039 *** & 0,038 *** \\
& (0,010) & (0,010) & (0,010) \\
25 a 34 años & 0,035 & 0,035 & 0,038 * \\
& (0,019) & (0,019) & (0,019) \\
35 a 44 años & 0,037 * & 0,036 * & 0,039 * \\
& (0,018) & (0,018) & (0,018) \\
45 a 54 años & 0,036 & 0,036 & 0,039 * \\
& (0,019) & (0,019) & (0,019) \\
55 a 64 años & 0,041 * & 0,040 * & 0,044 * \\
& (0,020) & (0,020) & (0,020) \\
Educación media & -0,014 & -0,014 & -0,017 \\
& (0,020) & (0,020) & (0,020) \\
Educación alta & -0,072 *** & -0,072 *** & -0,076 *** \\
& (0,021) & (0,021) & (0,020) \\
Decil de ingreso & -0,016 *** & -0,016 *** & -0,016 *** \\
& (0,002) & (0,002) & (0,002) \\
Asalariado temporal o sin contrato & 0,001 & 0,001 & 0,001 \\
& (0,015) & (0,015) & (0,015) \\
Independiente & 0,000 & 0,000 & 0,001 \\
& (0,015) & (0,015) & (0,015) \\
Trabajo en plataformas & -0,046 ** & -0,046 ** & -0,042 * \\
& (0,017) & (0,017) & (0,017) \\
Sector público & 0,004 & 0,004 & 0,001 \\
& (0,012) & (0,012) & (0,012) \\
\textbf{Nivel país} & & & \\
PIB per cápita (log) & -0,000 & -0,007 & -0,002 \\
& (0,021) & (0,019) & (0,017) \\
Gini & -0,041 & -0,044 & -0,022 \\
& (0,023) & (0,022) & (0,019) \\
Prestaciones por desempleo & -0,017 & & -0,023 \\
& (0,022) & & (0,020) \\
Políticas de mercado laboral & -0,050 * & -0,055 ** & -0,037 * \\
& (0,021) & (0,020) & (0,017) \\
Gasto social público & -0,054 ** & -0,062 *** & -0,036 * \\
& (0,021) & (0,018) & (0,017) \\
\textbf{Interacciones entre niveles} & & & \\
Amenaza tecnológica × Prestaciones por desempleo & & & 0,047 *** \\
& & & (0,011) \\
Log verosimilitud & -21267,088 & -21264,490 & -21222,845 \\
N & 19.731 & 19.731 & 19.731 \\
Países & 27 & 27 & 27 \\
Varianza del intercepto (país) & 0,005 & 0,005 & 0,005 \\
Varianza residual & 0,448 & 0,448 & 0,446 \\
Varianza de la pendiente de la amenaza tecnológica & & & 0,002 \\
Covarianza entre intercepto y pendiente (amenaza) & & & -0,002 \\

\end{longtable}

}

\emph{Nota}: variable dependiente: dirección de la compensación. El
tercer modelo incluye una pendiente aleatoria por país para la amenaza
tecnológica. Categorías de referencia: hombre, 18 a 24 años, educación
baja, asalariado con contrato indefinido, sin trabajo en plataformas y
no empleado en el sector público (asalariados del sector privado e
independientes). Todos los modelos controlan además por sector de
actividad. Errores estándar entre paréntesis. *** p \textless{} 0,001;
** p \textless{} 0,01; * p \textless{} 0,05. Los indicadores de nivel
país están estandarizados entre los 27 países. PIB per cápita en paridad
de poder adquisitivo, en logaritmo, y coeficiente de Gini, ambos del
Banco Mundial. Prestaciones por desempleo: porcentaje de personas
desempleadas que reciben prestaciones monetarias por desempleo (OIT,
indicador 1.3.1 de los ODS). Políticas de mercado laboral: gasto público
en políticas activas del mercado laboral como porcentaje del PIB,
dividido por la tasa de desempleo (OCDE, SOCX). Gasto social público:
gasto social público total como porcentaje del PIB (OCDE, SOCX).




\end{document}
